% ===================================================================
%  ka_makros.tex – Makros für Klassenarbeiten / Leistungskontrollen
%  Einbinden mit % ===================================================================
%  ka_makros.tex – Makros für Klassenarbeiten / Leistungskontrollen
%  Einbinden mit % ===================================================================
%  ka_makros.tex – Makros für Klassenarbeiten / Leistungskontrollen
%  Einbinden mit % ===================================================================
%  ka_makros.tex – Makros für Klassenarbeiten / Leistungskontrollen
%  Einbinden mit \input{ka_makros} nach \documentclass – oder nach
%  \input{header.tex} und \input{newcommands.tex} (universaler Stil).
%
%  \BE{2}            Bewertungseinheiten an eine (Teil-)Aufgabe hängen,
%                    halbe BE gehen auch: \BE{1.5}
%  \Kopf             Name/Klasse/Datum-Zeile + Angaben aus \Arbeit{...}
%  \Punktetabelle    Tabelle Aufgabe | BE | erreicht  (automatisch)
%  \Notenschluessel  Notentabelle nach \Grenzen (automatisch)
%  \loesungtrue      Lösungsteil mit ausgeben (\loesungfalse: ohne)
%
%  Alle Summen stehen erst nach dem ZWEITEN pdflatex-Lauf richtig da
%  (beim ersten Lauf erscheinen ??).
% ===================================================================
\makeatletter
% Pakete nur laden, wenn sie fehlen: so funktioniert die Datei allein
% (\documentclass + \input{ka_makros}) und nach der universalen header.tex
\@ifpackageloaded{fontenc}{}{\usepackage[T1]{fontenc}}
\@ifpackageloaded{lmodern}{}{\usepackage{lmodern}}
\@ifpackageloaded{babel}{}{\usepackage[ngerman]{babel}}
\@ifpackageloaded{geometry}{}{\usepackage[a4paper,margin=2cm]{geometry}}
\@ifpackageloaded{amsmath}{}{\usepackage{amsmath,amssymb}}
\@ifpackageloaded{xcolor}{}{\usepackage{xcolor}}
\@ifpackageloaded{enumitem}{}{\usepackage{enumitem}}
\@ifpackageloaded{multicol}{}{\usepackage{multicol}}
\@ifpackageloaded{booktabs}{}{\usepackage{array,booktabs}}
\@ifpackageloaded{etoolbox}{}{\usepackage{etoolbox}}
\@ifpackageloaded{tikz}{}{\usepackage{tikz}}
% Kein siunitx (kollidiert mit units aus der universalen header.tex):
% Zahlen mit Dezimalkomma über \KAzahl{1.5} -> 1,5
\newcommand\KAzahl[1]{\pgfmathprintnumber[use comma]{#1}}
\pagestyle{empty}                 % keine Kopf-/Fußzeilen
\setlength{\parindent}{0pt}


% --- Aufgabennummerierung (1), (2) … mit (a), (b) … ------------------
\setlist[enumerate,1]{label=(\arabic*), leftmargin=2.2em, itemsep=8pt}
\setlist[enumerate,2]{label=(\alph*), leftmargin=1.8em, itemsep=3pt}

% --- Schwierigkeitsstufen (falls nicht schon in newcommands.tex) ----
\providecommand\stufe[1]{%
  \ifcase#1\or
    \rule{1.3pt}{0.8ex}\,\textcolor{lightgray}{\rule{1.3pt}{1.3ex}\,\rule{1.3pt}{1.8ex}}\or
    \rule{1.3pt}{0.8ex}\,\rule{1.3pt}{1.3ex}\,\textcolor{lightgray}{\rule{1.3pt}{1.8ex}}\or
    \rule{1.3pt}{0.8ex}\,\rule{1.3pt}{1.3ex}\,\rule{1.3pt}{1.8ex}\fi\enspace}
\providecommand\leicht{\stufe{1}}
\providecommand\mittel{\stufe{2}}
\providecommand\schwer{\stufe{3}}

% --- Lösungsschalter --------------------------------------------------
\newif\ifloesung

% --- Bewertungseinheiten zählen ---------------------------------------
\gdef\BE@gesamt{0}
\gdef\BE@max{0}
\newcommand\BE[1]{%
  \edef\BE@nr{\arabic{enumi}}%
  \ifcsundef{BE@a\BE@nr}{\csgdef{BE@a\BE@nr}{0}}{}%
  \csxdef{BE@a\BE@nr}{\fpeval{\csuse{BE@a\BE@nr}+#1}}%
  \xdef\BE@gesamt{\fpeval{\BE@gesamt+#1}}%
  \ifnum\BE@nr>\BE@max\relax\xdef\BE@max{\BE@nr}\fi
  \unskip\nobreak\hfill\mbox{\enspace\textcolor{black!60}{\small[\KAzahl{#1}\,BE]}}}
% Am Ende in die .aux schreiben, damit der Kopf die Summe kennt
\AtEndDocument{%
  \immediate\write\@auxout{\string\gdef\string\BEsumme{\BE@gesamt}}%
  \immediate\write\@auxout{\string\gdef\string\BEanzahl{\BE@max}}%
  \@tempcnta=1\relax
  \loop\ifnum\@tempcnta>\BE@max\relax\else
    \immediate\write\@auxout{\string\csgdef{BEaufg\the\@tempcnta}{%
      \ifcsdef{BE@a\the\@tempcnta}{\csuse{BE@a\the\@tempcnta}}{0}}}%
    \advance\@tempcnta 1\relax\repeat}
\providecommand\BEsumme{0}  % erster Lauf: 0, ab dem zweiten Lauf korrekt
\providecommand\BEanzahl{0}

% --- Angaben zur Arbeit ------------------------------------------------
\newcommand\Arbeit[1]{\pgfkeys{/ka/.cd,#1}}
\pgfkeys{/ka/.cd,
  titel/.store in=\KAtitel,       titel=Klassenarbeit,
  fach/.store in=\KAfach,         fach=,
  klasse/.store in=\KAklasse,     klasse=,
  thema/.store in=\KAthema,       thema=,
  datum/.store in=\KAdatum,       datum=,
  zeit/.store in=\KAzeit,         zeit=,
  hilfsmittel/.store in=\KAhilfsmittel, hilfsmittel=,
  gruppe/.store in=\KAgruppe,     gruppe=}

\newcommand\Kopf{%
  {\sffamily\bfseries\Large \KAtitel\ \KAfach\ \KAklasse
   \ifx\KAgruppe\empty\else\hfill Gruppe \KAgruppe\fi}\par\smallskip
  {\sffamily\large \KAthema}\par\medskip
  \begin{tabular}{@{}p{0.52\linewidth}p{0.42\linewidth}@{}}
    Name: \hrulefill & Datum: \KAdatum \\[4pt]
    Arbeitszeit: \KAzeit & Hilfsmittel: \KAhilfsmittel \\[4pt]
    Erreichbar: \KAzahl{\BEsumme}\,BE & Erreicht: \rule{1.6cm}{0.4pt}\,BE
      \quad Note: \rule{1.2cm}{0.4pt}
  \end{tabular}\par
  \smallskip\hrule\medskip}

% --- Punktetabelle ---------------------------------------------------
\newcommand\Punktetabelle{%
  \gdef\PTkopf{}\gdef\PTbe{}\gdef\PTleer{}%
  \ifnum\BEanzahl>0
  \foreach \i in {1,...,\BEanzahl}{%
    \xappto\PTkopf{& (\i)}%
    \xappto\PTbe{& \noexpand\KAzahl{\csuse{BEaufg\i}}}%
    \gappto\PTleer{&}}%
  \fi
  \begin{tabular}{@{}l*{\BEanzahl}{>{\centering\arraybackslash}p{1cm}}|>{\centering\arraybackslash}p{1.3cm}@{}}
    Aufgabe \PTkopf & $\Sigma$ \\ \hline
    BE \PTbe & \KAzahl{\BEsumme} \\[3pt] \hline
    erreicht \PTleer & \rule{0pt}{16pt}\\ \hline
  \end{tabular}}

% --- Notenschlüssel ----------------------------------------------------
% \Grenzen{95,80,60,40,20}: Mindestprozent für die Noten 1 bis 5.
% Die BE-Grenze wird auf halbe BE AUFgerundet.
\newcommand\Grenzen[1]{\def\KAgrenzen{#1}}
\Grenzen{95,80,60,40,20}
\newcommand\abBE[1]{\fpeval{ceil(#1/100*\BEsumme*2)/2}}
\newcommand\Notenschluessel{%
  \edef\@g{\KAgrenzen}%
  \expandafter\KA@noten\@g\relax}
\def\KA@noten#1,#2,#3,#4,#5\relax{%
  \begingroup\small
  \begin{tabular}{@{}l*{6}{>{\centering\arraybackslash}p{1.45cm}}@{}}
    Note & 1 & 2 & 3 & 4 & 5 & 6 \\ \hline
    ab \% & \KAzahl{#1} & \KAzahl{#2} & \KAzahl{#3} & \KAzahl{#4} & \KAzahl{#5} & 0 \\
    BE & \KA@bereich{\abBE{#1}}{\BEsumme}
       & \KA@bereich{\abBE{#2}}{\fpeval{\abBE{#1}-0.5}}
       & \KA@bereich{\abBE{#3}}{\fpeval{\abBE{#2}-0.5}}
       & \KA@bereich{\abBE{#4}}{\fpeval{\abBE{#3}-0.5}}
       & \KA@bereich{\abBE{#5}}{\fpeval{\abBE{#4}-0.5}}
       & \KA@bereich{0}{\fpeval{\abBE{#5}-0.5}} \\[2pt] \hline
    Anzahl & & & & & & \rule{0pt}{14pt}\\
  \end{tabular}\endgroup}
\def\KA@bereich#1#2{\KAzahl{#1}\,--\,\KAzahl{#2}}

% --- Lösungsteil -------------------------------------------------------
% Lösungen parallel zum Aufgabenteil, \scriptsize, Ergebnisse in \boxed{}
\newenvironment{loesungsteil}{%
  \clearpage\scriptsize\raggedright
  {\sffamily\bfseries\normalsize Lösungen \KAtitel\ \KAfach\ \KAklasse}\par\smallskip
  \hrule\medskip}{}
\makeatother
 nach \documentclass – oder nach
%  \input{header.tex} und \input{newcommands.tex} (universaler Stil).
%
%  \BE{2}            Bewertungseinheiten an eine (Teil-)Aufgabe hängen,
%                    halbe BE gehen auch: \BE{1.5}
%  \Kopf             Name/Klasse/Datum-Zeile + Angaben aus \Arbeit{...}
%  \Punktetabelle    Tabelle Aufgabe | BE | erreicht  (automatisch)
%  \Notenschluessel  Notentabelle nach \Grenzen (automatisch)
%  \loesungtrue      Lösungsteil mit ausgeben (\loesungfalse: ohne)
%
%  Alle Summen stehen erst nach dem ZWEITEN pdflatex-Lauf richtig da
%  (beim ersten Lauf erscheinen ??).
% ===================================================================
\makeatletter
% Pakete nur laden, wenn sie fehlen: so funktioniert die Datei allein
% (\documentclass + % ===================================================================
%  ka_makros.tex – Makros für Klassenarbeiten / Leistungskontrollen
%  Einbinden mit \input{ka_makros} nach \documentclass – oder nach
%  \input{header.tex} und \input{newcommands.tex} (universaler Stil).
%
%  \BE{2}            Bewertungseinheiten an eine (Teil-)Aufgabe hängen,
%                    halbe BE gehen auch: \BE{1.5}
%  \Kopf             Name/Klasse/Datum-Zeile + Angaben aus \Arbeit{...}
%  \Punktetabelle    Tabelle Aufgabe | BE | erreicht  (automatisch)
%  \Notenschluessel  Notentabelle nach \Grenzen (automatisch)
%  \loesungtrue      Lösungsteil mit ausgeben (\loesungfalse: ohne)
%
%  Alle Summen stehen erst nach dem ZWEITEN pdflatex-Lauf richtig da
%  (beim ersten Lauf erscheinen ??).
% ===================================================================
\makeatletter
% Pakete nur laden, wenn sie fehlen: so funktioniert die Datei allein
% (\documentclass + \input{ka_makros}) und nach der universalen header.tex
\@ifpackageloaded{fontenc}{}{\usepackage[T1]{fontenc}}
\@ifpackageloaded{lmodern}{}{\usepackage{lmodern}}
\@ifpackageloaded{babel}{}{\usepackage[ngerman]{babel}}
\@ifpackageloaded{geometry}{}{\usepackage[a4paper,margin=2cm]{geometry}}
\@ifpackageloaded{amsmath}{}{\usepackage{amsmath,amssymb}}
\@ifpackageloaded{xcolor}{}{\usepackage{xcolor}}
\@ifpackageloaded{enumitem}{}{\usepackage{enumitem}}
\@ifpackageloaded{multicol}{}{\usepackage{multicol}}
\@ifpackageloaded{booktabs}{}{\usepackage{array,booktabs}}
\@ifpackageloaded{etoolbox}{}{\usepackage{etoolbox}}
\@ifpackageloaded{tikz}{}{\usepackage{tikz}}
% Kein siunitx (kollidiert mit units aus der universalen header.tex):
% Zahlen mit Dezimalkomma über \KAzahl{1.5} -> 1,5
\newcommand\KAzahl[1]{\pgfmathprintnumber[use comma]{#1}}
\pagestyle{empty}                 % keine Kopf-/Fußzeilen
\setlength{\parindent}{0pt}


% --- Aufgabennummerierung (1), (2) … mit (a), (b) … ------------------
\setlist[enumerate,1]{label=(\arabic*), leftmargin=2.2em, itemsep=8pt}
\setlist[enumerate,2]{label=(\alph*), leftmargin=1.8em, itemsep=3pt}

% --- Schwierigkeitsstufen (falls nicht schon in newcommands.tex) ----
\providecommand\stufe[1]{%
  \ifcase#1\or
    \rule{1.3pt}{0.8ex}\,\textcolor{lightgray}{\rule{1.3pt}{1.3ex}\,\rule{1.3pt}{1.8ex}}\or
    \rule{1.3pt}{0.8ex}\,\rule{1.3pt}{1.3ex}\,\textcolor{lightgray}{\rule{1.3pt}{1.8ex}}\or
    \rule{1.3pt}{0.8ex}\,\rule{1.3pt}{1.3ex}\,\rule{1.3pt}{1.8ex}\fi\enspace}
\providecommand\leicht{\stufe{1}}
\providecommand\mittel{\stufe{2}}
\providecommand\schwer{\stufe{3}}

% --- Lösungsschalter --------------------------------------------------
\newif\ifloesung

% --- Bewertungseinheiten zählen ---------------------------------------
\gdef\BE@gesamt{0}
\gdef\BE@max{0}
\newcommand\BE[1]{%
  \edef\BE@nr{\arabic{enumi}}%
  \ifcsundef{BE@a\BE@nr}{\csgdef{BE@a\BE@nr}{0}}{}%
  \csxdef{BE@a\BE@nr}{\fpeval{\csuse{BE@a\BE@nr}+#1}}%
  \xdef\BE@gesamt{\fpeval{\BE@gesamt+#1}}%
  \ifnum\BE@nr>\BE@max\relax\xdef\BE@max{\BE@nr}\fi
  \unskip\nobreak\hfill\mbox{\enspace\textcolor{black!60}{\small[\KAzahl{#1}\,BE]}}}
% Am Ende in die .aux schreiben, damit der Kopf die Summe kennt
\AtEndDocument{%
  \immediate\write\@auxout{\string\gdef\string\BEsumme{\BE@gesamt}}%
  \immediate\write\@auxout{\string\gdef\string\BEanzahl{\BE@max}}%
  \@tempcnta=1\relax
  \loop\ifnum\@tempcnta>\BE@max\relax\else
    \immediate\write\@auxout{\string\csgdef{BEaufg\the\@tempcnta}{%
      \ifcsdef{BE@a\the\@tempcnta}{\csuse{BE@a\the\@tempcnta}}{0}}}%
    \advance\@tempcnta 1\relax\repeat}
\providecommand\BEsumme{0}  % erster Lauf: 0, ab dem zweiten Lauf korrekt
\providecommand\BEanzahl{0}

% --- Angaben zur Arbeit ------------------------------------------------
\newcommand\Arbeit[1]{\pgfkeys{/ka/.cd,#1}}
\pgfkeys{/ka/.cd,
  titel/.store in=\KAtitel,       titel=Klassenarbeit,
  fach/.store in=\KAfach,         fach=,
  klasse/.store in=\KAklasse,     klasse=,
  thema/.store in=\KAthema,       thema=,
  datum/.store in=\KAdatum,       datum=,
  zeit/.store in=\KAzeit,         zeit=,
  hilfsmittel/.store in=\KAhilfsmittel, hilfsmittel=,
  gruppe/.store in=\KAgruppe,     gruppe=}

\newcommand\Kopf{%
  {\sffamily\bfseries\Large \KAtitel\ \KAfach\ \KAklasse
   \ifx\KAgruppe\empty\else\hfill Gruppe \KAgruppe\fi}\par\smallskip
  {\sffamily\large \KAthema}\par\medskip
  \begin{tabular}{@{}p{0.52\linewidth}p{0.42\linewidth}@{}}
    Name: \hrulefill & Datum: \KAdatum \\[4pt]
    Arbeitszeit: \KAzeit & Hilfsmittel: \KAhilfsmittel \\[4pt]
    Erreichbar: \KAzahl{\BEsumme}\,BE & Erreicht: \rule{1.6cm}{0.4pt}\,BE
      \quad Note: \rule{1.2cm}{0.4pt}
  \end{tabular}\par
  \smallskip\hrule\medskip}

% --- Punktetabelle ---------------------------------------------------
\newcommand\Punktetabelle{%
  \gdef\PTkopf{}\gdef\PTbe{}\gdef\PTleer{}%
  \ifnum\BEanzahl>0
  \foreach \i in {1,...,\BEanzahl}{%
    \xappto\PTkopf{& (\i)}%
    \xappto\PTbe{& \noexpand\KAzahl{\csuse{BEaufg\i}}}%
    \gappto\PTleer{&}}%
  \fi
  \begin{tabular}{@{}l*{\BEanzahl}{>{\centering\arraybackslash}p{1cm}}|>{\centering\arraybackslash}p{1.3cm}@{}}
    Aufgabe \PTkopf & $\Sigma$ \\ \hline
    BE \PTbe & \KAzahl{\BEsumme} \\[3pt] \hline
    erreicht \PTleer & \rule{0pt}{16pt}\\ \hline
  \end{tabular}}

% --- Notenschlüssel ----------------------------------------------------
% \Grenzen{95,80,60,40,20}: Mindestprozent für die Noten 1 bis 5.
% Die BE-Grenze wird auf halbe BE AUFgerundet.
\newcommand\Grenzen[1]{\def\KAgrenzen{#1}}
\Grenzen{95,80,60,40,20}
\newcommand\abBE[1]{\fpeval{ceil(#1/100*\BEsumme*2)/2}}
\newcommand\Notenschluessel{%
  \edef\@g{\KAgrenzen}%
  \expandafter\KA@noten\@g\relax}
\def\KA@noten#1,#2,#3,#4,#5\relax{%
  \begingroup\small
  \begin{tabular}{@{}l*{6}{>{\centering\arraybackslash}p{1.45cm}}@{}}
    Note & 1 & 2 & 3 & 4 & 5 & 6 \\ \hline
    ab \% & \KAzahl{#1} & \KAzahl{#2} & \KAzahl{#3} & \KAzahl{#4} & \KAzahl{#5} & 0 \\
    BE & \KA@bereich{\abBE{#1}}{\BEsumme}
       & \KA@bereich{\abBE{#2}}{\fpeval{\abBE{#1}-0.5}}
       & \KA@bereich{\abBE{#3}}{\fpeval{\abBE{#2}-0.5}}
       & \KA@bereich{\abBE{#4}}{\fpeval{\abBE{#3}-0.5}}
       & \KA@bereich{\abBE{#5}}{\fpeval{\abBE{#4}-0.5}}
       & \KA@bereich{0}{\fpeval{\abBE{#5}-0.5}} \\[2pt] \hline
    Anzahl & & & & & & \rule{0pt}{14pt}\\
  \end{tabular}\endgroup}
\def\KA@bereich#1#2{\KAzahl{#1}\,--\,\KAzahl{#2}}

% --- Lösungsteil -------------------------------------------------------
% Lösungen parallel zum Aufgabenteil, \scriptsize, Ergebnisse in \boxed{}
\newenvironment{loesungsteil}{%
  \clearpage\scriptsize\raggedright
  {\sffamily\bfseries\normalsize Lösungen \KAtitel\ \KAfach\ \KAklasse}\par\smallskip
  \hrule\medskip}{}
\makeatother
) und nach der universalen header.tex
\@ifpackageloaded{fontenc}{}{\usepackage[T1]{fontenc}}
\@ifpackageloaded{lmodern}{}{\usepackage{lmodern}}
\@ifpackageloaded{babel}{}{\usepackage[ngerman]{babel}}
\@ifpackageloaded{geometry}{}{\usepackage[a4paper,margin=2cm]{geometry}}
\@ifpackageloaded{amsmath}{}{\usepackage{amsmath,amssymb}}
\@ifpackageloaded{xcolor}{}{\usepackage{xcolor}}
\@ifpackageloaded{enumitem}{}{\usepackage{enumitem}}
\@ifpackageloaded{multicol}{}{\usepackage{multicol}}
\@ifpackageloaded{booktabs}{}{\usepackage{array,booktabs}}
\@ifpackageloaded{etoolbox}{}{\usepackage{etoolbox}}
\@ifpackageloaded{tikz}{}{\usepackage{tikz}}
% Kein siunitx (kollidiert mit units aus der universalen header.tex):
% Zahlen mit Dezimalkomma über \KAzahl{1.5} -> 1,5
\newcommand\KAzahl[1]{\pgfmathprintnumber[use comma]{#1}}
\pagestyle{empty}                 % keine Kopf-/Fußzeilen
\setlength{\parindent}{0pt}


% --- Aufgabennummerierung (1), (2) … mit (a), (b) … ------------------
\setlist[enumerate,1]{label=(\arabic*), leftmargin=2.2em, itemsep=8pt}
\setlist[enumerate,2]{label=(\alph*), leftmargin=1.8em, itemsep=3pt}

% --- Schwierigkeitsstufen (falls nicht schon in newcommands.tex) ----
\providecommand\stufe[1]{%
  \ifcase#1\or
    \rule{1.3pt}{0.8ex}\,\textcolor{lightgray}{\rule{1.3pt}{1.3ex}\,\rule{1.3pt}{1.8ex}}\or
    \rule{1.3pt}{0.8ex}\,\rule{1.3pt}{1.3ex}\,\textcolor{lightgray}{\rule{1.3pt}{1.8ex}}\or
    \rule{1.3pt}{0.8ex}\,\rule{1.3pt}{1.3ex}\,\rule{1.3pt}{1.8ex}\fi\enspace}
\providecommand\leicht{\stufe{1}}
\providecommand\mittel{\stufe{2}}
\providecommand\schwer{\stufe{3}}

% --- Lösungsschalter --------------------------------------------------
\newif\ifloesung

% --- Bewertungseinheiten zählen ---------------------------------------
\gdef\BE@gesamt{0}
\gdef\BE@max{0}
\newcommand\BE[1]{%
  \edef\BE@nr{\arabic{enumi}}%
  \ifcsundef{BE@a\BE@nr}{\csgdef{BE@a\BE@nr}{0}}{}%
  \csxdef{BE@a\BE@nr}{\fpeval{\csuse{BE@a\BE@nr}+#1}}%
  \xdef\BE@gesamt{\fpeval{\BE@gesamt+#1}}%
  \ifnum\BE@nr>\BE@max\relax\xdef\BE@max{\BE@nr}\fi
  \unskip\nobreak\hfill\mbox{\enspace\textcolor{black!60}{\small[\KAzahl{#1}\,BE]}}}
% Am Ende in die .aux schreiben, damit der Kopf die Summe kennt
\AtEndDocument{%
  \immediate\write\@auxout{\string\gdef\string\BEsumme{\BE@gesamt}}%
  \immediate\write\@auxout{\string\gdef\string\BEanzahl{\BE@max}}%
  \@tempcnta=1\relax
  \loop\ifnum\@tempcnta>\BE@max\relax\else
    \immediate\write\@auxout{\string\csgdef{BEaufg\the\@tempcnta}{%
      \ifcsdef{BE@a\the\@tempcnta}{\csuse{BE@a\the\@tempcnta}}{0}}}%
    \advance\@tempcnta 1\relax\repeat}
\providecommand\BEsumme{0}  % erster Lauf: 0, ab dem zweiten Lauf korrekt
\providecommand\BEanzahl{0}

% --- Angaben zur Arbeit ------------------------------------------------
\newcommand\Arbeit[1]{\pgfkeys{/ka/.cd,#1}}
\pgfkeys{/ka/.cd,
  titel/.store in=\KAtitel,       titel=Klassenarbeit,
  fach/.store in=\KAfach,         fach=,
  klasse/.store in=\KAklasse,     klasse=,
  thema/.store in=\KAthema,       thema=,
  datum/.store in=\KAdatum,       datum=,
  zeit/.store in=\KAzeit,         zeit=,
  hilfsmittel/.store in=\KAhilfsmittel, hilfsmittel=,
  gruppe/.store in=\KAgruppe,     gruppe=}

\newcommand\Kopf{%
  {\sffamily\bfseries\Large \KAtitel\ \KAfach\ \KAklasse
   \ifx\KAgruppe\empty\else\hfill Gruppe \KAgruppe\fi}\par\smallskip
  {\sffamily\large \KAthema}\par\medskip
  \begin{tabular}{@{}p{0.52\linewidth}p{0.42\linewidth}@{}}
    Name: \hrulefill & Datum: \KAdatum \\[4pt]
    Arbeitszeit: \KAzeit & Hilfsmittel: \KAhilfsmittel \\[4pt]
    Erreichbar: \KAzahl{\BEsumme}\,BE & Erreicht: \rule{1.6cm}{0.4pt}\,BE
      \quad Note: \rule{1.2cm}{0.4pt}
  \end{tabular}\par
  \smallskip\hrule\medskip}

% --- Punktetabelle ---------------------------------------------------
\newcommand\Punktetabelle{%
  \gdef\PTkopf{}\gdef\PTbe{}\gdef\PTleer{}%
  \ifnum\BEanzahl>0
  \foreach \i in {1,...,\BEanzahl}{%
    \xappto\PTkopf{& (\i)}%
    \xappto\PTbe{& \noexpand\KAzahl{\csuse{BEaufg\i}}}%
    \gappto\PTleer{&}}%
  \fi
  \begin{tabular}{@{}l*{\BEanzahl}{>{\centering\arraybackslash}p{1cm}}|>{\centering\arraybackslash}p{1.3cm}@{}}
    Aufgabe \PTkopf & $\Sigma$ \\ \hline
    BE \PTbe & \KAzahl{\BEsumme} \\[3pt] \hline
    erreicht \PTleer & \rule{0pt}{16pt}\\ \hline
  \end{tabular}}

% --- Notenschlüssel ----------------------------------------------------
% \Grenzen{95,80,60,40,20}: Mindestprozent für die Noten 1 bis 5.
% Die BE-Grenze wird auf halbe BE AUFgerundet.
\newcommand\Grenzen[1]{\def\KAgrenzen{#1}}
\Grenzen{95,80,60,40,20}
\newcommand\abBE[1]{\fpeval{ceil(#1/100*\BEsumme*2)/2}}
\newcommand\Notenschluessel{%
  \edef\@g{\KAgrenzen}%
  \expandafter\KA@noten\@g\relax}
\def\KA@noten#1,#2,#3,#4,#5\relax{%
  \begingroup\small
  \begin{tabular}{@{}l*{6}{>{\centering\arraybackslash}p{1.45cm}}@{}}
    Note & 1 & 2 & 3 & 4 & 5 & 6 \\ \hline
    ab \% & \KAzahl{#1} & \KAzahl{#2} & \KAzahl{#3} & \KAzahl{#4} & \KAzahl{#5} & 0 \\
    BE & \KA@bereich{\abBE{#1}}{\BEsumme}
       & \KA@bereich{\abBE{#2}}{\fpeval{\abBE{#1}-0.5}}
       & \KA@bereich{\abBE{#3}}{\fpeval{\abBE{#2}-0.5}}
       & \KA@bereich{\abBE{#4}}{\fpeval{\abBE{#3}-0.5}}
       & \KA@bereich{\abBE{#5}}{\fpeval{\abBE{#4}-0.5}}
       & \KA@bereich{0}{\fpeval{\abBE{#5}-0.5}} \\[2pt] \hline
    Anzahl & & & & & & \rule{0pt}{14pt}\\
  \end{tabular}\endgroup}
\def\KA@bereich#1#2{\KAzahl{#1}\,--\,\KAzahl{#2}}

% --- Lösungsteil -------------------------------------------------------
% Lösungen parallel zum Aufgabenteil, \scriptsize, Ergebnisse in \boxed{}
\newenvironment{loesungsteil}{%
  \clearpage\scriptsize\raggedright
  {\sffamily\bfseries\normalsize Lösungen \KAtitel\ \KAfach\ \KAklasse}\par\smallskip
  \hrule\medskip}{}
\makeatother
 nach \documentclass – oder nach
%  \input{header.tex} und \input{newcommands.tex} (universaler Stil).
%
%  \BE{2}            Bewertungseinheiten an eine (Teil-)Aufgabe hängen,
%                    halbe BE gehen auch: \BE{1.5}
%  \Kopf             Name/Klasse/Datum-Zeile + Angaben aus \Arbeit{...}
%  \Punktetabelle    Tabelle Aufgabe | BE | erreicht  (automatisch)
%  \Notenschluessel  Notentabelle nach \Grenzen (automatisch)
%  \loesungtrue      Lösungsteil mit ausgeben (\loesungfalse: ohne)
%
%  Alle Summen stehen erst nach dem ZWEITEN pdflatex-Lauf richtig da
%  (beim ersten Lauf erscheinen ??).
% ===================================================================
\makeatletter
% Pakete nur laden, wenn sie fehlen: so funktioniert die Datei allein
% (\documentclass + % ===================================================================
%  ka_makros.tex – Makros für Klassenarbeiten / Leistungskontrollen
%  Einbinden mit % ===================================================================
%  ka_makros.tex – Makros für Klassenarbeiten / Leistungskontrollen
%  Einbinden mit \input{ka_makros} nach \documentclass – oder nach
%  \input{header.tex} und \input{newcommands.tex} (universaler Stil).
%
%  \BE{2}            Bewertungseinheiten an eine (Teil-)Aufgabe hängen,
%                    halbe BE gehen auch: \BE{1.5}
%  \Kopf             Name/Klasse/Datum-Zeile + Angaben aus \Arbeit{...}
%  \Punktetabelle    Tabelle Aufgabe | BE | erreicht  (automatisch)
%  \Notenschluessel  Notentabelle nach \Grenzen (automatisch)
%  \loesungtrue      Lösungsteil mit ausgeben (\loesungfalse: ohne)
%
%  Alle Summen stehen erst nach dem ZWEITEN pdflatex-Lauf richtig da
%  (beim ersten Lauf erscheinen ??).
% ===================================================================
\makeatletter
% Pakete nur laden, wenn sie fehlen: so funktioniert die Datei allein
% (\documentclass + \input{ka_makros}) und nach der universalen header.tex
\@ifpackageloaded{fontenc}{}{\usepackage[T1]{fontenc}}
\@ifpackageloaded{lmodern}{}{\usepackage{lmodern}}
\@ifpackageloaded{babel}{}{\usepackage[ngerman]{babel}}
\@ifpackageloaded{geometry}{}{\usepackage[a4paper,margin=2cm]{geometry}}
\@ifpackageloaded{amsmath}{}{\usepackage{amsmath,amssymb}}
\@ifpackageloaded{xcolor}{}{\usepackage{xcolor}}
\@ifpackageloaded{enumitem}{}{\usepackage{enumitem}}
\@ifpackageloaded{multicol}{}{\usepackage{multicol}}
\@ifpackageloaded{booktabs}{}{\usepackage{array,booktabs}}
\@ifpackageloaded{etoolbox}{}{\usepackage{etoolbox}}
\@ifpackageloaded{tikz}{}{\usepackage{tikz}}
% Kein siunitx (kollidiert mit units aus der universalen header.tex):
% Zahlen mit Dezimalkomma über \KAzahl{1.5} -> 1,5
\newcommand\KAzahl[1]{\pgfmathprintnumber[use comma]{#1}}
\pagestyle{empty}                 % keine Kopf-/Fußzeilen
\setlength{\parindent}{0pt}


% --- Aufgabennummerierung (1), (2) … mit (a), (b) … ------------------
\setlist[enumerate,1]{label=(\arabic*), leftmargin=2.2em, itemsep=8pt}
\setlist[enumerate,2]{label=(\alph*), leftmargin=1.8em, itemsep=3pt}

% --- Schwierigkeitsstufen (falls nicht schon in newcommands.tex) ----
\providecommand\stufe[1]{%
  \ifcase#1\or
    \rule{1.3pt}{0.8ex}\,\textcolor{lightgray}{\rule{1.3pt}{1.3ex}\,\rule{1.3pt}{1.8ex}}\or
    \rule{1.3pt}{0.8ex}\,\rule{1.3pt}{1.3ex}\,\textcolor{lightgray}{\rule{1.3pt}{1.8ex}}\or
    \rule{1.3pt}{0.8ex}\,\rule{1.3pt}{1.3ex}\,\rule{1.3pt}{1.8ex}\fi\enspace}
\providecommand\leicht{\stufe{1}}
\providecommand\mittel{\stufe{2}}
\providecommand\schwer{\stufe{3}}

% --- Lösungsschalter --------------------------------------------------
\newif\ifloesung

% --- Bewertungseinheiten zählen ---------------------------------------
\gdef\BE@gesamt{0}
\gdef\BE@max{0}
\newcommand\BE[1]{%
  \edef\BE@nr{\arabic{enumi}}%
  \ifcsundef{BE@a\BE@nr}{\csgdef{BE@a\BE@nr}{0}}{}%
  \csxdef{BE@a\BE@nr}{\fpeval{\csuse{BE@a\BE@nr}+#1}}%
  \xdef\BE@gesamt{\fpeval{\BE@gesamt+#1}}%
  \ifnum\BE@nr>\BE@max\relax\xdef\BE@max{\BE@nr}\fi
  \unskip\nobreak\hfill\mbox{\enspace\textcolor{black!60}{\small[\KAzahl{#1}\,BE]}}}
% Am Ende in die .aux schreiben, damit der Kopf die Summe kennt
\AtEndDocument{%
  \immediate\write\@auxout{\string\gdef\string\BEsumme{\BE@gesamt}}%
  \immediate\write\@auxout{\string\gdef\string\BEanzahl{\BE@max}}%
  \@tempcnta=1\relax
  \loop\ifnum\@tempcnta>\BE@max\relax\else
    \immediate\write\@auxout{\string\csgdef{BEaufg\the\@tempcnta}{%
      \ifcsdef{BE@a\the\@tempcnta}{\csuse{BE@a\the\@tempcnta}}{0}}}%
    \advance\@tempcnta 1\relax\repeat}
\providecommand\BEsumme{0}  % erster Lauf: 0, ab dem zweiten Lauf korrekt
\providecommand\BEanzahl{0}

% --- Angaben zur Arbeit ------------------------------------------------
\newcommand\Arbeit[1]{\pgfkeys{/ka/.cd,#1}}
\pgfkeys{/ka/.cd,
  titel/.store in=\KAtitel,       titel=Klassenarbeit,
  fach/.store in=\KAfach,         fach=,
  klasse/.store in=\KAklasse,     klasse=,
  thema/.store in=\KAthema,       thema=,
  datum/.store in=\KAdatum,       datum=,
  zeit/.store in=\KAzeit,         zeit=,
  hilfsmittel/.store in=\KAhilfsmittel, hilfsmittel=,
  gruppe/.store in=\KAgruppe,     gruppe=}

\newcommand\Kopf{%
  {\sffamily\bfseries\Large \KAtitel\ \KAfach\ \KAklasse
   \ifx\KAgruppe\empty\else\hfill Gruppe \KAgruppe\fi}\par\smallskip
  {\sffamily\large \KAthema}\par\medskip
  \begin{tabular}{@{}p{0.52\linewidth}p{0.42\linewidth}@{}}
    Name: \hrulefill & Datum: \KAdatum \\[4pt]
    Arbeitszeit: \KAzeit & Hilfsmittel: \KAhilfsmittel \\[4pt]
    Erreichbar: \KAzahl{\BEsumme}\,BE & Erreicht: \rule{1.6cm}{0.4pt}\,BE
      \quad Note: \rule{1.2cm}{0.4pt}
  \end{tabular}\par
  \smallskip\hrule\medskip}

% --- Punktetabelle ---------------------------------------------------
\newcommand\Punktetabelle{%
  \gdef\PTkopf{}\gdef\PTbe{}\gdef\PTleer{}%
  \ifnum\BEanzahl>0
  \foreach \i in {1,...,\BEanzahl}{%
    \xappto\PTkopf{& (\i)}%
    \xappto\PTbe{& \noexpand\KAzahl{\csuse{BEaufg\i}}}%
    \gappto\PTleer{&}}%
  \fi
  \begin{tabular}{@{}l*{\BEanzahl}{>{\centering\arraybackslash}p{1cm}}|>{\centering\arraybackslash}p{1.3cm}@{}}
    Aufgabe \PTkopf & $\Sigma$ \\ \hline
    BE \PTbe & \KAzahl{\BEsumme} \\[3pt] \hline
    erreicht \PTleer & \rule{0pt}{16pt}\\ \hline
  \end{tabular}}

% --- Notenschlüssel ----------------------------------------------------
% \Grenzen{95,80,60,40,20}: Mindestprozent für die Noten 1 bis 5.
% Die BE-Grenze wird auf halbe BE AUFgerundet.
\newcommand\Grenzen[1]{\def\KAgrenzen{#1}}
\Grenzen{95,80,60,40,20}
\newcommand\abBE[1]{\fpeval{ceil(#1/100*\BEsumme*2)/2}}
\newcommand\Notenschluessel{%
  \edef\@g{\KAgrenzen}%
  \expandafter\KA@noten\@g\relax}
\def\KA@noten#1,#2,#3,#4,#5\relax{%
  \begingroup\small
  \begin{tabular}{@{}l*{6}{>{\centering\arraybackslash}p{1.45cm}}@{}}
    Note & 1 & 2 & 3 & 4 & 5 & 6 \\ \hline
    ab \% & \KAzahl{#1} & \KAzahl{#2} & \KAzahl{#3} & \KAzahl{#4} & \KAzahl{#5} & 0 \\
    BE & \KA@bereich{\abBE{#1}}{\BEsumme}
       & \KA@bereich{\abBE{#2}}{\fpeval{\abBE{#1}-0.5}}
       & \KA@bereich{\abBE{#3}}{\fpeval{\abBE{#2}-0.5}}
       & \KA@bereich{\abBE{#4}}{\fpeval{\abBE{#3}-0.5}}
       & \KA@bereich{\abBE{#5}}{\fpeval{\abBE{#4}-0.5}}
       & \KA@bereich{0}{\fpeval{\abBE{#5}-0.5}} \\[2pt] \hline
    Anzahl & & & & & & \rule{0pt}{14pt}\\
  \end{tabular}\endgroup}
\def\KA@bereich#1#2{\KAzahl{#1}\,--\,\KAzahl{#2}}

% --- Lösungsteil -------------------------------------------------------
% Lösungen parallel zum Aufgabenteil, \scriptsize, Ergebnisse in \boxed{}
\newenvironment{loesungsteil}{%
  \clearpage\scriptsize\raggedright
  {\sffamily\bfseries\normalsize Lösungen \KAtitel\ \KAfach\ \KAklasse}\par\smallskip
  \hrule\medskip}{}
\makeatother
 nach \documentclass – oder nach
%  \input{header.tex} und \input{newcommands.tex} (universaler Stil).
%
%  \BE{2}            Bewertungseinheiten an eine (Teil-)Aufgabe hängen,
%                    halbe BE gehen auch: \BE{1.5}
%  \Kopf             Name/Klasse/Datum-Zeile + Angaben aus \Arbeit{...}
%  \Punktetabelle    Tabelle Aufgabe | BE | erreicht  (automatisch)
%  \Notenschluessel  Notentabelle nach \Grenzen (automatisch)
%  \loesungtrue      Lösungsteil mit ausgeben (\loesungfalse: ohne)
%
%  Alle Summen stehen erst nach dem ZWEITEN pdflatex-Lauf richtig da
%  (beim ersten Lauf erscheinen ??).
% ===================================================================
\makeatletter
% Pakete nur laden, wenn sie fehlen: so funktioniert die Datei allein
% (\documentclass + % ===================================================================
%  ka_makros.tex – Makros für Klassenarbeiten / Leistungskontrollen
%  Einbinden mit \input{ka_makros} nach \documentclass – oder nach
%  \input{header.tex} und \input{newcommands.tex} (universaler Stil).
%
%  \BE{2}            Bewertungseinheiten an eine (Teil-)Aufgabe hängen,
%                    halbe BE gehen auch: \BE{1.5}
%  \Kopf             Name/Klasse/Datum-Zeile + Angaben aus \Arbeit{...}
%  \Punktetabelle    Tabelle Aufgabe | BE | erreicht  (automatisch)
%  \Notenschluessel  Notentabelle nach \Grenzen (automatisch)
%  \loesungtrue      Lösungsteil mit ausgeben (\loesungfalse: ohne)
%
%  Alle Summen stehen erst nach dem ZWEITEN pdflatex-Lauf richtig da
%  (beim ersten Lauf erscheinen ??).
% ===================================================================
\makeatletter
% Pakete nur laden, wenn sie fehlen: so funktioniert die Datei allein
% (\documentclass + \input{ka_makros}) und nach der universalen header.tex
\@ifpackageloaded{fontenc}{}{\usepackage[T1]{fontenc}}
\@ifpackageloaded{lmodern}{}{\usepackage{lmodern}}
\@ifpackageloaded{babel}{}{\usepackage[ngerman]{babel}}
\@ifpackageloaded{geometry}{}{\usepackage[a4paper,margin=2cm]{geometry}}
\@ifpackageloaded{amsmath}{}{\usepackage{amsmath,amssymb}}
\@ifpackageloaded{xcolor}{}{\usepackage{xcolor}}
\@ifpackageloaded{enumitem}{}{\usepackage{enumitem}}
\@ifpackageloaded{multicol}{}{\usepackage{multicol}}
\@ifpackageloaded{booktabs}{}{\usepackage{array,booktabs}}
\@ifpackageloaded{etoolbox}{}{\usepackage{etoolbox}}
\@ifpackageloaded{tikz}{}{\usepackage{tikz}}
% Kein siunitx (kollidiert mit units aus der universalen header.tex):
% Zahlen mit Dezimalkomma über \KAzahl{1.5} -> 1,5
\newcommand\KAzahl[1]{\pgfmathprintnumber[use comma]{#1}}
\pagestyle{empty}                 % keine Kopf-/Fußzeilen
\setlength{\parindent}{0pt}


% --- Aufgabennummerierung (1), (2) … mit (a), (b) … ------------------
\setlist[enumerate,1]{label=(\arabic*), leftmargin=2.2em, itemsep=8pt}
\setlist[enumerate,2]{label=(\alph*), leftmargin=1.8em, itemsep=3pt}

% --- Schwierigkeitsstufen (falls nicht schon in newcommands.tex) ----
\providecommand\stufe[1]{%
  \ifcase#1\or
    \rule{1.3pt}{0.8ex}\,\textcolor{lightgray}{\rule{1.3pt}{1.3ex}\,\rule{1.3pt}{1.8ex}}\or
    \rule{1.3pt}{0.8ex}\,\rule{1.3pt}{1.3ex}\,\textcolor{lightgray}{\rule{1.3pt}{1.8ex}}\or
    \rule{1.3pt}{0.8ex}\,\rule{1.3pt}{1.3ex}\,\rule{1.3pt}{1.8ex}\fi\enspace}
\providecommand\leicht{\stufe{1}}
\providecommand\mittel{\stufe{2}}
\providecommand\schwer{\stufe{3}}

% --- Lösungsschalter --------------------------------------------------
\newif\ifloesung

% --- Bewertungseinheiten zählen ---------------------------------------
\gdef\BE@gesamt{0}
\gdef\BE@max{0}
\newcommand\BE[1]{%
  \edef\BE@nr{\arabic{enumi}}%
  \ifcsundef{BE@a\BE@nr}{\csgdef{BE@a\BE@nr}{0}}{}%
  \csxdef{BE@a\BE@nr}{\fpeval{\csuse{BE@a\BE@nr}+#1}}%
  \xdef\BE@gesamt{\fpeval{\BE@gesamt+#1}}%
  \ifnum\BE@nr>\BE@max\relax\xdef\BE@max{\BE@nr}\fi
  \unskip\nobreak\hfill\mbox{\enspace\textcolor{black!60}{\small[\KAzahl{#1}\,BE]}}}
% Am Ende in die .aux schreiben, damit der Kopf die Summe kennt
\AtEndDocument{%
  \immediate\write\@auxout{\string\gdef\string\BEsumme{\BE@gesamt}}%
  \immediate\write\@auxout{\string\gdef\string\BEanzahl{\BE@max}}%
  \@tempcnta=1\relax
  \loop\ifnum\@tempcnta>\BE@max\relax\else
    \immediate\write\@auxout{\string\csgdef{BEaufg\the\@tempcnta}{%
      \ifcsdef{BE@a\the\@tempcnta}{\csuse{BE@a\the\@tempcnta}}{0}}}%
    \advance\@tempcnta 1\relax\repeat}
\providecommand\BEsumme{0}  % erster Lauf: 0, ab dem zweiten Lauf korrekt
\providecommand\BEanzahl{0}

% --- Angaben zur Arbeit ------------------------------------------------
\newcommand\Arbeit[1]{\pgfkeys{/ka/.cd,#1}}
\pgfkeys{/ka/.cd,
  titel/.store in=\KAtitel,       titel=Klassenarbeit,
  fach/.store in=\KAfach,         fach=,
  klasse/.store in=\KAklasse,     klasse=,
  thema/.store in=\KAthema,       thema=,
  datum/.store in=\KAdatum,       datum=,
  zeit/.store in=\KAzeit,         zeit=,
  hilfsmittel/.store in=\KAhilfsmittel, hilfsmittel=,
  gruppe/.store in=\KAgruppe,     gruppe=}

\newcommand\Kopf{%
  {\sffamily\bfseries\Large \KAtitel\ \KAfach\ \KAklasse
   \ifx\KAgruppe\empty\else\hfill Gruppe \KAgruppe\fi}\par\smallskip
  {\sffamily\large \KAthema}\par\medskip
  \begin{tabular}{@{}p{0.52\linewidth}p{0.42\linewidth}@{}}
    Name: \hrulefill & Datum: \KAdatum \\[4pt]
    Arbeitszeit: \KAzeit & Hilfsmittel: \KAhilfsmittel \\[4pt]
    Erreichbar: \KAzahl{\BEsumme}\,BE & Erreicht: \rule{1.6cm}{0.4pt}\,BE
      \quad Note: \rule{1.2cm}{0.4pt}
  \end{tabular}\par
  \smallskip\hrule\medskip}

% --- Punktetabelle ---------------------------------------------------
\newcommand\Punktetabelle{%
  \gdef\PTkopf{}\gdef\PTbe{}\gdef\PTleer{}%
  \ifnum\BEanzahl>0
  \foreach \i in {1,...,\BEanzahl}{%
    \xappto\PTkopf{& (\i)}%
    \xappto\PTbe{& \noexpand\KAzahl{\csuse{BEaufg\i}}}%
    \gappto\PTleer{&}}%
  \fi
  \begin{tabular}{@{}l*{\BEanzahl}{>{\centering\arraybackslash}p{1cm}}|>{\centering\arraybackslash}p{1.3cm}@{}}
    Aufgabe \PTkopf & $\Sigma$ \\ \hline
    BE \PTbe & \KAzahl{\BEsumme} \\[3pt] \hline
    erreicht \PTleer & \rule{0pt}{16pt}\\ \hline
  \end{tabular}}

% --- Notenschlüssel ----------------------------------------------------
% \Grenzen{95,80,60,40,20}: Mindestprozent für die Noten 1 bis 5.
% Die BE-Grenze wird auf halbe BE AUFgerundet.
\newcommand\Grenzen[1]{\def\KAgrenzen{#1}}
\Grenzen{95,80,60,40,20}
\newcommand\abBE[1]{\fpeval{ceil(#1/100*\BEsumme*2)/2}}
\newcommand\Notenschluessel{%
  \edef\@g{\KAgrenzen}%
  \expandafter\KA@noten\@g\relax}
\def\KA@noten#1,#2,#3,#4,#5\relax{%
  \begingroup\small
  \begin{tabular}{@{}l*{6}{>{\centering\arraybackslash}p{1.45cm}}@{}}
    Note & 1 & 2 & 3 & 4 & 5 & 6 \\ \hline
    ab \% & \KAzahl{#1} & \KAzahl{#2} & \KAzahl{#3} & \KAzahl{#4} & \KAzahl{#5} & 0 \\
    BE & \KA@bereich{\abBE{#1}}{\BEsumme}
       & \KA@bereich{\abBE{#2}}{\fpeval{\abBE{#1}-0.5}}
       & \KA@bereich{\abBE{#3}}{\fpeval{\abBE{#2}-0.5}}
       & \KA@bereich{\abBE{#4}}{\fpeval{\abBE{#3}-0.5}}
       & \KA@bereich{\abBE{#5}}{\fpeval{\abBE{#4}-0.5}}
       & \KA@bereich{0}{\fpeval{\abBE{#5}-0.5}} \\[2pt] \hline
    Anzahl & & & & & & \rule{0pt}{14pt}\\
  \end{tabular}\endgroup}
\def\KA@bereich#1#2{\KAzahl{#1}\,--\,\KAzahl{#2}}

% --- Lösungsteil -------------------------------------------------------
% Lösungen parallel zum Aufgabenteil, \scriptsize, Ergebnisse in \boxed{}
\newenvironment{loesungsteil}{%
  \clearpage\scriptsize\raggedright
  {\sffamily\bfseries\normalsize Lösungen \KAtitel\ \KAfach\ \KAklasse}\par\smallskip
  \hrule\medskip}{}
\makeatother
) und nach der universalen header.tex
\@ifpackageloaded{fontenc}{}{\usepackage[T1]{fontenc}}
\@ifpackageloaded{lmodern}{}{\usepackage{lmodern}}
\@ifpackageloaded{babel}{}{\usepackage[ngerman]{babel}}
\@ifpackageloaded{geometry}{}{\usepackage[a4paper,margin=2cm]{geometry}}
\@ifpackageloaded{amsmath}{}{\usepackage{amsmath,amssymb}}
\@ifpackageloaded{xcolor}{}{\usepackage{xcolor}}
\@ifpackageloaded{enumitem}{}{\usepackage{enumitem}}
\@ifpackageloaded{multicol}{}{\usepackage{multicol}}
\@ifpackageloaded{booktabs}{}{\usepackage{array,booktabs}}
\@ifpackageloaded{etoolbox}{}{\usepackage{etoolbox}}
\@ifpackageloaded{tikz}{}{\usepackage{tikz}}
% Kein siunitx (kollidiert mit units aus der universalen header.tex):
% Zahlen mit Dezimalkomma über \KAzahl{1.5} -> 1,5
\newcommand\KAzahl[1]{\pgfmathprintnumber[use comma]{#1}}
\pagestyle{empty}                 % keine Kopf-/Fußzeilen
\setlength{\parindent}{0pt}


% --- Aufgabennummerierung (1), (2) … mit (a), (b) … ------------------
\setlist[enumerate,1]{label=(\arabic*), leftmargin=2.2em, itemsep=8pt}
\setlist[enumerate,2]{label=(\alph*), leftmargin=1.8em, itemsep=3pt}

% --- Schwierigkeitsstufen (falls nicht schon in newcommands.tex) ----
\providecommand\stufe[1]{%
  \ifcase#1\or
    \rule{1.3pt}{0.8ex}\,\textcolor{lightgray}{\rule{1.3pt}{1.3ex}\,\rule{1.3pt}{1.8ex}}\or
    \rule{1.3pt}{0.8ex}\,\rule{1.3pt}{1.3ex}\,\textcolor{lightgray}{\rule{1.3pt}{1.8ex}}\or
    \rule{1.3pt}{0.8ex}\,\rule{1.3pt}{1.3ex}\,\rule{1.3pt}{1.8ex}\fi\enspace}
\providecommand\leicht{\stufe{1}}
\providecommand\mittel{\stufe{2}}
\providecommand\schwer{\stufe{3}}

% --- Lösungsschalter --------------------------------------------------
\newif\ifloesung

% --- Bewertungseinheiten zählen ---------------------------------------
\gdef\BE@gesamt{0}
\gdef\BE@max{0}
\newcommand\BE[1]{%
  \edef\BE@nr{\arabic{enumi}}%
  \ifcsundef{BE@a\BE@nr}{\csgdef{BE@a\BE@nr}{0}}{}%
  \csxdef{BE@a\BE@nr}{\fpeval{\csuse{BE@a\BE@nr}+#1}}%
  \xdef\BE@gesamt{\fpeval{\BE@gesamt+#1}}%
  \ifnum\BE@nr>\BE@max\relax\xdef\BE@max{\BE@nr}\fi
  \unskip\nobreak\hfill\mbox{\enspace\textcolor{black!60}{\small[\KAzahl{#1}\,BE]}}}
% Am Ende in die .aux schreiben, damit der Kopf die Summe kennt
\AtEndDocument{%
  \immediate\write\@auxout{\string\gdef\string\BEsumme{\BE@gesamt}}%
  \immediate\write\@auxout{\string\gdef\string\BEanzahl{\BE@max}}%
  \@tempcnta=1\relax
  \loop\ifnum\@tempcnta>\BE@max\relax\else
    \immediate\write\@auxout{\string\csgdef{BEaufg\the\@tempcnta}{%
      \ifcsdef{BE@a\the\@tempcnta}{\csuse{BE@a\the\@tempcnta}}{0}}}%
    \advance\@tempcnta 1\relax\repeat}
\providecommand\BEsumme{0}  % erster Lauf: 0, ab dem zweiten Lauf korrekt
\providecommand\BEanzahl{0}

% --- Angaben zur Arbeit ------------------------------------------------
\newcommand\Arbeit[1]{\pgfkeys{/ka/.cd,#1}}
\pgfkeys{/ka/.cd,
  titel/.store in=\KAtitel,       titel=Klassenarbeit,
  fach/.store in=\KAfach,         fach=,
  klasse/.store in=\KAklasse,     klasse=,
  thema/.store in=\KAthema,       thema=,
  datum/.store in=\KAdatum,       datum=,
  zeit/.store in=\KAzeit,         zeit=,
  hilfsmittel/.store in=\KAhilfsmittel, hilfsmittel=,
  gruppe/.store in=\KAgruppe,     gruppe=}

\newcommand\Kopf{%
  {\sffamily\bfseries\Large \KAtitel\ \KAfach\ \KAklasse
   \ifx\KAgruppe\empty\else\hfill Gruppe \KAgruppe\fi}\par\smallskip
  {\sffamily\large \KAthema}\par\medskip
  \begin{tabular}{@{}p{0.52\linewidth}p{0.42\linewidth}@{}}
    Name: \hrulefill & Datum: \KAdatum \\[4pt]
    Arbeitszeit: \KAzeit & Hilfsmittel: \KAhilfsmittel \\[4pt]
    Erreichbar: \KAzahl{\BEsumme}\,BE & Erreicht: \rule{1.6cm}{0.4pt}\,BE
      \quad Note: \rule{1.2cm}{0.4pt}
  \end{tabular}\par
  \smallskip\hrule\medskip}

% --- Punktetabelle ---------------------------------------------------
\newcommand\Punktetabelle{%
  \gdef\PTkopf{}\gdef\PTbe{}\gdef\PTleer{}%
  \ifnum\BEanzahl>0
  \foreach \i in {1,...,\BEanzahl}{%
    \xappto\PTkopf{& (\i)}%
    \xappto\PTbe{& \noexpand\KAzahl{\csuse{BEaufg\i}}}%
    \gappto\PTleer{&}}%
  \fi
  \begin{tabular}{@{}l*{\BEanzahl}{>{\centering\arraybackslash}p{1cm}}|>{\centering\arraybackslash}p{1.3cm}@{}}
    Aufgabe \PTkopf & $\Sigma$ \\ \hline
    BE \PTbe & \KAzahl{\BEsumme} \\[3pt] \hline
    erreicht \PTleer & \rule{0pt}{16pt}\\ \hline
  \end{tabular}}

% --- Notenschlüssel ----------------------------------------------------
% \Grenzen{95,80,60,40,20}: Mindestprozent für die Noten 1 bis 5.
% Die BE-Grenze wird auf halbe BE AUFgerundet.
\newcommand\Grenzen[1]{\def\KAgrenzen{#1}}
\Grenzen{95,80,60,40,20}
\newcommand\abBE[1]{\fpeval{ceil(#1/100*\BEsumme*2)/2}}
\newcommand\Notenschluessel{%
  \edef\@g{\KAgrenzen}%
  \expandafter\KA@noten\@g\relax}
\def\KA@noten#1,#2,#3,#4,#5\relax{%
  \begingroup\small
  \begin{tabular}{@{}l*{6}{>{\centering\arraybackslash}p{1.45cm}}@{}}
    Note & 1 & 2 & 3 & 4 & 5 & 6 \\ \hline
    ab \% & \KAzahl{#1} & \KAzahl{#2} & \KAzahl{#3} & \KAzahl{#4} & \KAzahl{#5} & 0 \\
    BE & \KA@bereich{\abBE{#1}}{\BEsumme}
       & \KA@bereich{\abBE{#2}}{\fpeval{\abBE{#1}-0.5}}
       & \KA@bereich{\abBE{#3}}{\fpeval{\abBE{#2}-0.5}}
       & \KA@bereich{\abBE{#4}}{\fpeval{\abBE{#3}-0.5}}
       & \KA@bereich{\abBE{#5}}{\fpeval{\abBE{#4}-0.5}}
       & \KA@bereich{0}{\fpeval{\abBE{#5}-0.5}} \\[2pt] \hline
    Anzahl & & & & & & \rule{0pt}{14pt}\\
  \end{tabular}\endgroup}
\def\KA@bereich#1#2{\KAzahl{#1}\,--\,\KAzahl{#2}}

% --- Lösungsteil -------------------------------------------------------
% Lösungen parallel zum Aufgabenteil, \scriptsize, Ergebnisse in \boxed{}
\newenvironment{loesungsteil}{%
  \clearpage\scriptsize\raggedright
  {\sffamily\bfseries\normalsize Lösungen \KAtitel\ \KAfach\ \KAklasse}\par\smallskip
  \hrule\medskip}{}
\makeatother
) und nach der universalen header.tex
\@ifpackageloaded{fontenc}{}{\usepackage[T1]{fontenc}}
\@ifpackageloaded{lmodern}{}{\usepackage{lmodern}}
\@ifpackageloaded{babel}{}{\usepackage[ngerman]{babel}}
\@ifpackageloaded{geometry}{}{\usepackage[a4paper,margin=2cm]{geometry}}
\@ifpackageloaded{amsmath}{}{\usepackage{amsmath,amssymb}}
\@ifpackageloaded{xcolor}{}{\usepackage{xcolor}}
\@ifpackageloaded{enumitem}{}{\usepackage{enumitem}}
\@ifpackageloaded{multicol}{}{\usepackage{multicol}}
\@ifpackageloaded{booktabs}{}{\usepackage{array,booktabs}}
\@ifpackageloaded{etoolbox}{}{\usepackage{etoolbox}}
\@ifpackageloaded{tikz}{}{\usepackage{tikz}}
% Kein siunitx (kollidiert mit units aus der universalen header.tex):
% Zahlen mit Dezimalkomma über \KAzahl{1.5} -> 1,5
\newcommand\KAzahl[1]{\pgfmathprintnumber[use comma]{#1}}
\pagestyle{empty}                 % keine Kopf-/Fußzeilen
\setlength{\parindent}{0pt}


% --- Aufgabennummerierung (1), (2) … mit (a), (b) … ------------------
\setlist[enumerate,1]{label=(\arabic*), leftmargin=2.2em, itemsep=8pt}
\setlist[enumerate,2]{label=(\alph*), leftmargin=1.8em, itemsep=3pt}

% --- Schwierigkeitsstufen (falls nicht schon in newcommands.tex) ----
\providecommand\stufe[1]{%
  \ifcase#1\or
    \rule{1.3pt}{0.8ex}\,\textcolor{lightgray}{\rule{1.3pt}{1.3ex}\,\rule{1.3pt}{1.8ex}}\or
    \rule{1.3pt}{0.8ex}\,\rule{1.3pt}{1.3ex}\,\textcolor{lightgray}{\rule{1.3pt}{1.8ex}}\or
    \rule{1.3pt}{0.8ex}\,\rule{1.3pt}{1.3ex}\,\rule{1.3pt}{1.8ex}\fi\enspace}
\providecommand\leicht{\stufe{1}}
\providecommand\mittel{\stufe{2}}
\providecommand\schwer{\stufe{3}}

% --- Lösungsschalter --------------------------------------------------
\newif\ifloesung

% --- Bewertungseinheiten zählen ---------------------------------------
\gdef\BE@gesamt{0}
\gdef\BE@max{0}
\newcommand\BE[1]{%
  \edef\BE@nr{\arabic{enumi}}%
  \ifcsundef{BE@a\BE@nr}{\csgdef{BE@a\BE@nr}{0}}{}%
  \csxdef{BE@a\BE@nr}{\fpeval{\csuse{BE@a\BE@nr}+#1}}%
  \xdef\BE@gesamt{\fpeval{\BE@gesamt+#1}}%
  \ifnum\BE@nr>\BE@max\relax\xdef\BE@max{\BE@nr}\fi
  \unskip\nobreak\hfill\mbox{\enspace\textcolor{black!60}{\small[\KAzahl{#1}\,BE]}}}
% Am Ende in die .aux schreiben, damit der Kopf die Summe kennt
\AtEndDocument{%
  \immediate\write\@auxout{\string\gdef\string\BEsumme{\BE@gesamt}}%
  \immediate\write\@auxout{\string\gdef\string\BEanzahl{\BE@max}}%
  \@tempcnta=1\relax
  \loop\ifnum\@tempcnta>\BE@max\relax\else
    \immediate\write\@auxout{\string\csgdef{BEaufg\the\@tempcnta}{%
      \ifcsdef{BE@a\the\@tempcnta}{\csuse{BE@a\the\@tempcnta}}{0}}}%
    \advance\@tempcnta 1\relax\repeat}
\providecommand\BEsumme{0}  % erster Lauf: 0, ab dem zweiten Lauf korrekt
\providecommand\BEanzahl{0}

% --- Angaben zur Arbeit ------------------------------------------------
\newcommand\Arbeit[1]{\pgfkeys{/ka/.cd,#1}}
\pgfkeys{/ka/.cd,
  titel/.store in=\KAtitel,       titel=Klassenarbeit,
  fach/.store in=\KAfach,         fach=,
  klasse/.store in=\KAklasse,     klasse=,
  thema/.store in=\KAthema,       thema=,
  datum/.store in=\KAdatum,       datum=,
  zeit/.store in=\KAzeit,         zeit=,
  hilfsmittel/.store in=\KAhilfsmittel, hilfsmittel=,
  gruppe/.store in=\KAgruppe,     gruppe=}

\newcommand\Kopf{%
  {\sffamily\bfseries\Large \KAtitel\ \KAfach\ \KAklasse
   \ifx\KAgruppe\empty\else\hfill Gruppe \KAgruppe\fi}\par\smallskip
  {\sffamily\large \KAthema}\par\medskip
  \begin{tabular}{@{}p{0.52\linewidth}p{0.42\linewidth}@{}}
    Name: \hrulefill & Datum: \KAdatum \\[4pt]
    Arbeitszeit: \KAzeit & Hilfsmittel: \KAhilfsmittel \\[4pt]
    Erreichbar: \KAzahl{\BEsumme}\,BE & Erreicht: \rule{1.6cm}{0.4pt}\,BE
      \quad Note: \rule{1.2cm}{0.4pt}
  \end{tabular}\par
  \smallskip\hrule\medskip}

% --- Punktetabelle ---------------------------------------------------
\newcommand\Punktetabelle{%
  \gdef\PTkopf{}\gdef\PTbe{}\gdef\PTleer{}%
  \ifnum\BEanzahl>0
  \foreach \i in {1,...,\BEanzahl}{%
    \xappto\PTkopf{& (\i)}%
    \xappto\PTbe{& \noexpand\KAzahl{\csuse{BEaufg\i}}}%
    \gappto\PTleer{&}}%
  \fi
  \begin{tabular}{@{}l*{\BEanzahl}{>{\centering\arraybackslash}p{1cm}}|>{\centering\arraybackslash}p{1.3cm}@{}}
    Aufgabe \PTkopf & $\Sigma$ \\ \hline
    BE \PTbe & \KAzahl{\BEsumme} \\[3pt] \hline
    erreicht \PTleer & \rule{0pt}{16pt}\\ \hline
  \end{tabular}}

% --- Notenschlüssel ----------------------------------------------------
% \Grenzen{95,80,60,40,20}: Mindestprozent für die Noten 1 bis 5.
% Die BE-Grenze wird auf halbe BE AUFgerundet.
\newcommand\Grenzen[1]{\def\KAgrenzen{#1}}
\Grenzen{95,80,60,40,20}
\newcommand\abBE[1]{\fpeval{ceil(#1/100*\BEsumme*2)/2}}
\newcommand\Notenschluessel{%
  \edef\@g{\KAgrenzen}%
  \expandafter\KA@noten\@g\relax}
\def\KA@noten#1,#2,#3,#4,#5\relax{%
  \begingroup\small
  \begin{tabular}{@{}l*{6}{>{\centering\arraybackslash}p{1.45cm}}@{}}
    Note & 1 & 2 & 3 & 4 & 5 & 6 \\ \hline
    ab \% & \KAzahl{#1} & \KAzahl{#2} & \KAzahl{#3} & \KAzahl{#4} & \KAzahl{#5} & 0 \\
    BE & \KA@bereich{\abBE{#1}}{\BEsumme}
       & \KA@bereich{\abBE{#2}}{\fpeval{\abBE{#1}-0.5}}
       & \KA@bereich{\abBE{#3}}{\fpeval{\abBE{#2}-0.5}}
       & \KA@bereich{\abBE{#4}}{\fpeval{\abBE{#3}-0.5}}
       & \KA@bereich{\abBE{#5}}{\fpeval{\abBE{#4}-0.5}}
       & \KA@bereich{0}{\fpeval{\abBE{#5}-0.5}} \\[2pt] \hline
    Anzahl & & & & & & \rule{0pt}{14pt}\\
  \end{tabular}\endgroup}
\def\KA@bereich#1#2{\KAzahl{#1}\,--\,\KAzahl{#2}}

% --- Lösungsteil -------------------------------------------------------
% Lösungen parallel zum Aufgabenteil, \scriptsize, Ergebnisse in \boxed{}
\newenvironment{loesungsteil}{%
  \clearpage\scriptsize\raggedright
  {\sffamily\bfseries\normalsize Lösungen \KAtitel\ \KAfach\ \KAklasse}\par\smallskip
  \hrule\medskip}{}
\makeatother
 nach \documentclass – oder nach
%  \input{header.tex} und \input{newcommands.tex} (universaler Stil).
%
%  \BE{2}            Bewertungseinheiten an eine (Teil-)Aufgabe hängen,
%                    halbe BE gehen auch: \BE{1.5}
%  \Kopf             Name/Klasse/Datum-Zeile + Angaben aus \Arbeit{...}
%  \Punktetabelle    Tabelle Aufgabe | BE | erreicht  (automatisch)
%  \Notenschluessel  Notentabelle nach \Grenzen (automatisch)
%  \loesungtrue      Lösungsteil mit ausgeben (\loesungfalse: ohne)
%
%  Alle Summen stehen erst nach dem ZWEITEN pdflatex-Lauf richtig da
%  (beim ersten Lauf erscheinen ??).
% ===================================================================
\makeatletter
% Pakete nur laden, wenn sie fehlen: so funktioniert die Datei allein
% (\documentclass + % ===================================================================
%  ka_makros.tex – Makros für Klassenarbeiten / Leistungskontrollen
%  Einbinden mit % ===================================================================
%  ka_makros.tex – Makros für Klassenarbeiten / Leistungskontrollen
%  Einbinden mit % ===================================================================
%  ka_makros.tex – Makros für Klassenarbeiten / Leistungskontrollen
%  Einbinden mit \input{ka_makros} nach \documentclass – oder nach
%  \input{header.tex} und \input{newcommands.tex} (universaler Stil).
%
%  \BE{2}            Bewertungseinheiten an eine (Teil-)Aufgabe hängen,
%                    halbe BE gehen auch: \BE{1.5}
%  \Kopf             Name/Klasse/Datum-Zeile + Angaben aus \Arbeit{...}
%  \Punktetabelle    Tabelle Aufgabe | BE | erreicht  (automatisch)
%  \Notenschluessel  Notentabelle nach \Grenzen (automatisch)
%  \loesungtrue      Lösungsteil mit ausgeben (\loesungfalse: ohne)
%
%  Alle Summen stehen erst nach dem ZWEITEN pdflatex-Lauf richtig da
%  (beim ersten Lauf erscheinen ??).
% ===================================================================
\makeatletter
% Pakete nur laden, wenn sie fehlen: so funktioniert die Datei allein
% (\documentclass + \input{ka_makros}) und nach der universalen header.tex
\@ifpackageloaded{fontenc}{}{\usepackage[T1]{fontenc}}
\@ifpackageloaded{lmodern}{}{\usepackage{lmodern}}
\@ifpackageloaded{babel}{}{\usepackage[ngerman]{babel}}
\@ifpackageloaded{geometry}{}{\usepackage[a4paper,margin=2cm]{geometry}}
\@ifpackageloaded{amsmath}{}{\usepackage{amsmath,amssymb}}
\@ifpackageloaded{xcolor}{}{\usepackage{xcolor}}
\@ifpackageloaded{enumitem}{}{\usepackage{enumitem}}
\@ifpackageloaded{multicol}{}{\usepackage{multicol}}
\@ifpackageloaded{booktabs}{}{\usepackage{array,booktabs}}
\@ifpackageloaded{etoolbox}{}{\usepackage{etoolbox}}
\@ifpackageloaded{tikz}{}{\usepackage{tikz}}
% Kein siunitx (kollidiert mit units aus der universalen header.tex):
% Zahlen mit Dezimalkomma über \KAzahl{1.5} -> 1,5
\newcommand\KAzahl[1]{\pgfmathprintnumber[use comma]{#1}}
\pagestyle{empty}                 % keine Kopf-/Fußzeilen
\setlength{\parindent}{0pt}


% --- Aufgabennummerierung (1), (2) … mit (a), (b) … ------------------
\setlist[enumerate,1]{label=(\arabic*), leftmargin=2.2em, itemsep=8pt}
\setlist[enumerate,2]{label=(\alph*), leftmargin=1.8em, itemsep=3pt}

% --- Schwierigkeitsstufen (falls nicht schon in newcommands.tex) ----
\providecommand\stufe[1]{%
  \ifcase#1\or
    \rule{1.3pt}{0.8ex}\,\textcolor{lightgray}{\rule{1.3pt}{1.3ex}\,\rule{1.3pt}{1.8ex}}\or
    \rule{1.3pt}{0.8ex}\,\rule{1.3pt}{1.3ex}\,\textcolor{lightgray}{\rule{1.3pt}{1.8ex}}\or
    \rule{1.3pt}{0.8ex}\,\rule{1.3pt}{1.3ex}\,\rule{1.3pt}{1.8ex}\fi\enspace}
\providecommand\leicht{\stufe{1}}
\providecommand\mittel{\stufe{2}}
\providecommand\schwer{\stufe{3}}

% --- Lösungsschalter --------------------------------------------------
\newif\ifloesung

% --- Bewertungseinheiten zählen ---------------------------------------
\gdef\BE@gesamt{0}
\gdef\BE@max{0}
\newcommand\BE[1]{%
  \edef\BE@nr{\arabic{enumi}}%
  \ifcsundef{BE@a\BE@nr}{\csgdef{BE@a\BE@nr}{0}}{}%
  \csxdef{BE@a\BE@nr}{\fpeval{\csuse{BE@a\BE@nr}+#1}}%
  \xdef\BE@gesamt{\fpeval{\BE@gesamt+#1}}%
  \ifnum\BE@nr>\BE@max\relax\xdef\BE@max{\BE@nr}\fi
  \unskip\nobreak\hfill\mbox{\enspace\textcolor{black!60}{\small[\KAzahl{#1}\,BE]}}}
% Am Ende in die .aux schreiben, damit der Kopf die Summe kennt
\AtEndDocument{%
  \immediate\write\@auxout{\string\gdef\string\BEsumme{\BE@gesamt}}%
  \immediate\write\@auxout{\string\gdef\string\BEanzahl{\BE@max}}%
  \@tempcnta=1\relax
  \loop\ifnum\@tempcnta>\BE@max\relax\else
    \immediate\write\@auxout{\string\csgdef{BEaufg\the\@tempcnta}{%
      \ifcsdef{BE@a\the\@tempcnta}{\csuse{BE@a\the\@tempcnta}}{0}}}%
    \advance\@tempcnta 1\relax\repeat}
\providecommand\BEsumme{0}  % erster Lauf: 0, ab dem zweiten Lauf korrekt
\providecommand\BEanzahl{0}

% --- Angaben zur Arbeit ------------------------------------------------
\newcommand\Arbeit[1]{\pgfkeys{/ka/.cd,#1}}
\pgfkeys{/ka/.cd,
  titel/.store in=\KAtitel,       titel=Klassenarbeit,
  fach/.store in=\KAfach,         fach=,
  klasse/.store in=\KAklasse,     klasse=,
  thema/.store in=\KAthema,       thema=,
  datum/.store in=\KAdatum,       datum=,
  zeit/.store in=\KAzeit,         zeit=,
  hilfsmittel/.store in=\KAhilfsmittel, hilfsmittel=,
  gruppe/.store in=\KAgruppe,     gruppe=}

\newcommand\Kopf{%
  {\sffamily\bfseries\Large \KAtitel\ \KAfach\ \KAklasse
   \ifx\KAgruppe\empty\else\hfill Gruppe \KAgruppe\fi}\par\smallskip
  {\sffamily\large \KAthema}\par\medskip
  \begin{tabular}{@{}p{0.52\linewidth}p{0.42\linewidth}@{}}
    Name: \hrulefill & Datum: \KAdatum \\[4pt]
    Arbeitszeit: \KAzeit & Hilfsmittel: \KAhilfsmittel \\[4pt]
    Erreichbar: \KAzahl{\BEsumme}\,BE & Erreicht: \rule{1.6cm}{0.4pt}\,BE
      \quad Note: \rule{1.2cm}{0.4pt}
  \end{tabular}\par
  \smallskip\hrule\medskip}

% --- Punktetabelle ---------------------------------------------------
\newcommand\Punktetabelle{%
  \gdef\PTkopf{}\gdef\PTbe{}\gdef\PTleer{}%
  \ifnum\BEanzahl>0
  \foreach \i in {1,...,\BEanzahl}{%
    \xappto\PTkopf{& (\i)}%
    \xappto\PTbe{& \noexpand\KAzahl{\csuse{BEaufg\i}}}%
    \gappto\PTleer{&}}%
  \fi
  \begin{tabular}{@{}l*{\BEanzahl}{>{\centering\arraybackslash}p{1cm}}|>{\centering\arraybackslash}p{1.3cm}@{}}
    Aufgabe \PTkopf & $\Sigma$ \\ \hline
    BE \PTbe & \KAzahl{\BEsumme} \\[3pt] \hline
    erreicht \PTleer & \rule{0pt}{16pt}\\ \hline
  \end{tabular}}

% --- Notenschlüssel ----------------------------------------------------
% \Grenzen{95,80,60,40,20}: Mindestprozent für die Noten 1 bis 5.
% Die BE-Grenze wird auf halbe BE AUFgerundet.
\newcommand\Grenzen[1]{\def\KAgrenzen{#1}}
\Grenzen{95,80,60,40,20}
\newcommand\abBE[1]{\fpeval{ceil(#1/100*\BEsumme*2)/2}}
\newcommand\Notenschluessel{%
  \edef\@g{\KAgrenzen}%
  \expandafter\KA@noten\@g\relax}
\def\KA@noten#1,#2,#3,#4,#5\relax{%
  \begingroup\small
  \begin{tabular}{@{}l*{6}{>{\centering\arraybackslash}p{1.45cm}}@{}}
    Note & 1 & 2 & 3 & 4 & 5 & 6 \\ \hline
    ab \% & \KAzahl{#1} & \KAzahl{#2} & \KAzahl{#3} & \KAzahl{#4} & \KAzahl{#5} & 0 \\
    BE & \KA@bereich{\abBE{#1}}{\BEsumme}
       & \KA@bereich{\abBE{#2}}{\fpeval{\abBE{#1}-0.5}}
       & \KA@bereich{\abBE{#3}}{\fpeval{\abBE{#2}-0.5}}
       & \KA@bereich{\abBE{#4}}{\fpeval{\abBE{#3}-0.5}}
       & \KA@bereich{\abBE{#5}}{\fpeval{\abBE{#4}-0.5}}
       & \KA@bereich{0}{\fpeval{\abBE{#5}-0.5}} \\[2pt] \hline
    Anzahl & & & & & & \rule{0pt}{14pt}\\
  \end{tabular}\endgroup}
\def\KA@bereich#1#2{\KAzahl{#1}\,--\,\KAzahl{#2}}

% --- Lösungsteil -------------------------------------------------------
% Lösungen parallel zum Aufgabenteil, \scriptsize, Ergebnisse in \boxed{}
\newenvironment{loesungsteil}{%
  \clearpage\scriptsize\raggedright
  {\sffamily\bfseries\normalsize Lösungen \KAtitel\ \KAfach\ \KAklasse}\par\smallskip
  \hrule\medskip}{}
\makeatother
 nach \documentclass – oder nach
%  \input{header.tex} und \input{newcommands.tex} (universaler Stil).
%
%  \BE{2}            Bewertungseinheiten an eine (Teil-)Aufgabe hängen,
%                    halbe BE gehen auch: \BE{1.5}
%  \Kopf             Name/Klasse/Datum-Zeile + Angaben aus \Arbeit{...}
%  \Punktetabelle    Tabelle Aufgabe | BE | erreicht  (automatisch)
%  \Notenschluessel  Notentabelle nach \Grenzen (automatisch)
%  \loesungtrue      Lösungsteil mit ausgeben (\loesungfalse: ohne)
%
%  Alle Summen stehen erst nach dem ZWEITEN pdflatex-Lauf richtig da
%  (beim ersten Lauf erscheinen ??).
% ===================================================================
\makeatletter
% Pakete nur laden, wenn sie fehlen: so funktioniert die Datei allein
% (\documentclass + % ===================================================================
%  ka_makros.tex – Makros für Klassenarbeiten / Leistungskontrollen
%  Einbinden mit \input{ka_makros} nach \documentclass – oder nach
%  \input{header.tex} und \input{newcommands.tex} (universaler Stil).
%
%  \BE{2}            Bewertungseinheiten an eine (Teil-)Aufgabe hängen,
%                    halbe BE gehen auch: \BE{1.5}
%  \Kopf             Name/Klasse/Datum-Zeile + Angaben aus \Arbeit{...}
%  \Punktetabelle    Tabelle Aufgabe | BE | erreicht  (automatisch)
%  \Notenschluessel  Notentabelle nach \Grenzen (automatisch)
%  \loesungtrue      Lösungsteil mit ausgeben (\loesungfalse: ohne)
%
%  Alle Summen stehen erst nach dem ZWEITEN pdflatex-Lauf richtig da
%  (beim ersten Lauf erscheinen ??).
% ===================================================================
\makeatletter
% Pakete nur laden, wenn sie fehlen: so funktioniert die Datei allein
% (\documentclass + \input{ka_makros}) und nach der universalen header.tex
\@ifpackageloaded{fontenc}{}{\usepackage[T1]{fontenc}}
\@ifpackageloaded{lmodern}{}{\usepackage{lmodern}}
\@ifpackageloaded{babel}{}{\usepackage[ngerman]{babel}}
\@ifpackageloaded{geometry}{}{\usepackage[a4paper,margin=2cm]{geometry}}
\@ifpackageloaded{amsmath}{}{\usepackage{amsmath,amssymb}}
\@ifpackageloaded{xcolor}{}{\usepackage{xcolor}}
\@ifpackageloaded{enumitem}{}{\usepackage{enumitem}}
\@ifpackageloaded{multicol}{}{\usepackage{multicol}}
\@ifpackageloaded{booktabs}{}{\usepackage{array,booktabs}}
\@ifpackageloaded{etoolbox}{}{\usepackage{etoolbox}}
\@ifpackageloaded{tikz}{}{\usepackage{tikz}}
% Kein siunitx (kollidiert mit units aus der universalen header.tex):
% Zahlen mit Dezimalkomma über \KAzahl{1.5} -> 1,5
\newcommand\KAzahl[1]{\pgfmathprintnumber[use comma]{#1}}
\pagestyle{empty}                 % keine Kopf-/Fußzeilen
\setlength{\parindent}{0pt}


% --- Aufgabennummerierung (1), (2) … mit (a), (b) … ------------------
\setlist[enumerate,1]{label=(\arabic*), leftmargin=2.2em, itemsep=8pt}
\setlist[enumerate,2]{label=(\alph*), leftmargin=1.8em, itemsep=3pt}

% --- Schwierigkeitsstufen (falls nicht schon in newcommands.tex) ----
\providecommand\stufe[1]{%
  \ifcase#1\or
    \rule{1.3pt}{0.8ex}\,\textcolor{lightgray}{\rule{1.3pt}{1.3ex}\,\rule{1.3pt}{1.8ex}}\or
    \rule{1.3pt}{0.8ex}\,\rule{1.3pt}{1.3ex}\,\textcolor{lightgray}{\rule{1.3pt}{1.8ex}}\or
    \rule{1.3pt}{0.8ex}\,\rule{1.3pt}{1.3ex}\,\rule{1.3pt}{1.8ex}\fi\enspace}
\providecommand\leicht{\stufe{1}}
\providecommand\mittel{\stufe{2}}
\providecommand\schwer{\stufe{3}}

% --- Lösungsschalter --------------------------------------------------
\newif\ifloesung

% --- Bewertungseinheiten zählen ---------------------------------------
\gdef\BE@gesamt{0}
\gdef\BE@max{0}
\newcommand\BE[1]{%
  \edef\BE@nr{\arabic{enumi}}%
  \ifcsundef{BE@a\BE@nr}{\csgdef{BE@a\BE@nr}{0}}{}%
  \csxdef{BE@a\BE@nr}{\fpeval{\csuse{BE@a\BE@nr}+#1}}%
  \xdef\BE@gesamt{\fpeval{\BE@gesamt+#1}}%
  \ifnum\BE@nr>\BE@max\relax\xdef\BE@max{\BE@nr}\fi
  \unskip\nobreak\hfill\mbox{\enspace\textcolor{black!60}{\small[\KAzahl{#1}\,BE]}}}
% Am Ende in die .aux schreiben, damit der Kopf die Summe kennt
\AtEndDocument{%
  \immediate\write\@auxout{\string\gdef\string\BEsumme{\BE@gesamt}}%
  \immediate\write\@auxout{\string\gdef\string\BEanzahl{\BE@max}}%
  \@tempcnta=1\relax
  \loop\ifnum\@tempcnta>\BE@max\relax\else
    \immediate\write\@auxout{\string\csgdef{BEaufg\the\@tempcnta}{%
      \ifcsdef{BE@a\the\@tempcnta}{\csuse{BE@a\the\@tempcnta}}{0}}}%
    \advance\@tempcnta 1\relax\repeat}
\providecommand\BEsumme{0}  % erster Lauf: 0, ab dem zweiten Lauf korrekt
\providecommand\BEanzahl{0}

% --- Angaben zur Arbeit ------------------------------------------------
\newcommand\Arbeit[1]{\pgfkeys{/ka/.cd,#1}}
\pgfkeys{/ka/.cd,
  titel/.store in=\KAtitel,       titel=Klassenarbeit,
  fach/.store in=\KAfach,         fach=,
  klasse/.store in=\KAklasse,     klasse=,
  thema/.store in=\KAthema,       thema=,
  datum/.store in=\KAdatum,       datum=,
  zeit/.store in=\KAzeit,         zeit=,
  hilfsmittel/.store in=\KAhilfsmittel, hilfsmittel=,
  gruppe/.store in=\KAgruppe,     gruppe=}

\newcommand\Kopf{%
  {\sffamily\bfseries\Large \KAtitel\ \KAfach\ \KAklasse
   \ifx\KAgruppe\empty\else\hfill Gruppe \KAgruppe\fi}\par\smallskip
  {\sffamily\large \KAthema}\par\medskip
  \begin{tabular}{@{}p{0.52\linewidth}p{0.42\linewidth}@{}}
    Name: \hrulefill & Datum: \KAdatum \\[4pt]
    Arbeitszeit: \KAzeit & Hilfsmittel: \KAhilfsmittel \\[4pt]
    Erreichbar: \KAzahl{\BEsumme}\,BE & Erreicht: \rule{1.6cm}{0.4pt}\,BE
      \quad Note: \rule{1.2cm}{0.4pt}
  \end{tabular}\par
  \smallskip\hrule\medskip}

% --- Punktetabelle ---------------------------------------------------
\newcommand\Punktetabelle{%
  \gdef\PTkopf{}\gdef\PTbe{}\gdef\PTleer{}%
  \ifnum\BEanzahl>0
  \foreach \i in {1,...,\BEanzahl}{%
    \xappto\PTkopf{& (\i)}%
    \xappto\PTbe{& \noexpand\KAzahl{\csuse{BEaufg\i}}}%
    \gappto\PTleer{&}}%
  \fi
  \begin{tabular}{@{}l*{\BEanzahl}{>{\centering\arraybackslash}p{1cm}}|>{\centering\arraybackslash}p{1.3cm}@{}}
    Aufgabe \PTkopf & $\Sigma$ \\ \hline
    BE \PTbe & \KAzahl{\BEsumme} \\[3pt] \hline
    erreicht \PTleer & \rule{0pt}{16pt}\\ \hline
  \end{tabular}}

% --- Notenschlüssel ----------------------------------------------------
% \Grenzen{95,80,60,40,20}: Mindestprozent für die Noten 1 bis 5.
% Die BE-Grenze wird auf halbe BE AUFgerundet.
\newcommand\Grenzen[1]{\def\KAgrenzen{#1}}
\Grenzen{95,80,60,40,20}
\newcommand\abBE[1]{\fpeval{ceil(#1/100*\BEsumme*2)/2}}
\newcommand\Notenschluessel{%
  \edef\@g{\KAgrenzen}%
  \expandafter\KA@noten\@g\relax}
\def\KA@noten#1,#2,#3,#4,#5\relax{%
  \begingroup\small
  \begin{tabular}{@{}l*{6}{>{\centering\arraybackslash}p{1.45cm}}@{}}
    Note & 1 & 2 & 3 & 4 & 5 & 6 \\ \hline
    ab \% & \KAzahl{#1} & \KAzahl{#2} & \KAzahl{#3} & \KAzahl{#4} & \KAzahl{#5} & 0 \\
    BE & \KA@bereich{\abBE{#1}}{\BEsumme}
       & \KA@bereich{\abBE{#2}}{\fpeval{\abBE{#1}-0.5}}
       & \KA@bereich{\abBE{#3}}{\fpeval{\abBE{#2}-0.5}}
       & \KA@bereich{\abBE{#4}}{\fpeval{\abBE{#3}-0.5}}
       & \KA@bereich{\abBE{#5}}{\fpeval{\abBE{#4}-0.5}}
       & \KA@bereich{0}{\fpeval{\abBE{#5}-0.5}} \\[2pt] \hline
    Anzahl & & & & & & \rule{0pt}{14pt}\\
  \end{tabular}\endgroup}
\def\KA@bereich#1#2{\KAzahl{#1}\,--\,\KAzahl{#2}}

% --- Lösungsteil -------------------------------------------------------
% Lösungen parallel zum Aufgabenteil, \scriptsize, Ergebnisse in \boxed{}
\newenvironment{loesungsteil}{%
  \clearpage\scriptsize\raggedright
  {\sffamily\bfseries\normalsize Lösungen \KAtitel\ \KAfach\ \KAklasse}\par\smallskip
  \hrule\medskip}{}
\makeatother
) und nach der universalen header.tex
\@ifpackageloaded{fontenc}{}{\usepackage[T1]{fontenc}}
\@ifpackageloaded{lmodern}{}{\usepackage{lmodern}}
\@ifpackageloaded{babel}{}{\usepackage[ngerman]{babel}}
\@ifpackageloaded{geometry}{}{\usepackage[a4paper,margin=2cm]{geometry}}
\@ifpackageloaded{amsmath}{}{\usepackage{amsmath,amssymb}}
\@ifpackageloaded{xcolor}{}{\usepackage{xcolor}}
\@ifpackageloaded{enumitem}{}{\usepackage{enumitem}}
\@ifpackageloaded{multicol}{}{\usepackage{multicol}}
\@ifpackageloaded{booktabs}{}{\usepackage{array,booktabs}}
\@ifpackageloaded{etoolbox}{}{\usepackage{etoolbox}}
\@ifpackageloaded{tikz}{}{\usepackage{tikz}}
% Kein siunitx (kollidiert mit units aus der universalen header.tex):
% Zahlen mit Dezimalkomma über \KAzahl{1.5} -> 1,5
\newcommand\KAzahl[1]{\pgfmathprintnumber[use comma]{#1}}
\pagestyle{empty}                 % keine Kopf-/Fußzeilen
\setlength{\parindent}{0pt}


% --- Aufgabennummerierung (1), (2) … mit (a), (b) … ------------------
\setlist[enumerate,1]{label=(\arabic*), leftmargin=2.2em, itemsep=8pt}
\setlist[enumerate,2]{label=(\alph*), leftmargin=1.8em, itemsep=3pt}

% --- Schwierigkeitsstufen (falls nicht schon in newcommands.tex) ----
\providecommand\stufe[1]{%
  \ifcase#1\or
    \rule{1.3pt}{0.8ex}\,\textcolor{lightgray}{\rule{1.3pt}{1.3ex}\,\rule{1.3pt}{1.8ex}}\or
    \rule{1.3pt}{0.8ex}\,\rule{1.3pt}{1.3ex}\,\textcolor{lightgray}{\rule{1.3pt}{1.8ex}}\or
    \rule{1.3pt}{0.8ex}\,\rule{1.3pt}{1.3ex}\,\rule{1.3pt}{1.8ex}\fi\enspace}
\providecommand\leicht{\stufe{1}}
\providecommand\mittel{\stufe{2}}
\providecommand\schwer{\stufe{3}}

% --- Lösungsschalter --------------------------------------------------
\newif\ifloesung

% --- Bewertungseinheiten zählen ---------------------------------------
\gdef\BE@gesamt{0}
\gdef\BE@max{0}
\newcommand\BE[1]{%
  \edef\BE@nr{\arabic{enumi}}%
  \ifcsundef{BE@a\BE@nr}{\csgdef{BE@a\BE@nr}{0}}{}%
  \csxdef{BE@a\BE@nr}{\fpeval{\csuse{BE@a\BE@nr}+#1}}%
  \xdef\BE@gesamt{\fpeval{\BE@gesamt+#1}}%
  \ifnum\BE@nr>\BE@max\relax\xdef\BE@max{\BE@nr}\fi
  \unskip\nobreak\hfill\mbox{\enspace\textcolor{black!60}{\small[\KAzahl{#1}\,BE]}}}
% Am Ende in die .aux schreiben, damit der Kopf die Summe kennt
\AtEndDocument{%
  \immediate\write\@auxout{\string\gdef\string\BEsumme{\BE@gesamt}}%
  \immediate\write\@auxout{\string\gdef\string\BEanzahl{\BE@max}}%
  \@tempcnta=1\relax
  \loop\ifnum\@tempcnta>\BE@max\relax\else
    \immediate\write\@auxout{\string\csgdef{BEaufg\the\@tempcnta}{%
      \ifcsdef{BE@a\the\@tempcnta}{\csuse{BE@a\the\@tempcnta}}{0}}}%
    \advance\@tempcnta 1\relax\repeat}
\providecommand\BEsumme{0}  % erster Lauf: 0, ab dem zweiten Lauf korrekt
\providecommand\BEanzahl{0}

% --- Angaben zur Arbeit ------------------------------------------------
\newcommand\Arbeit[1]{\pgfkeys{/ka/.cd,#1}}
\pgfkeys{/ka/.cd,
  titel/.store in=\KAtitel,       titel=Klassenarbeit,
  fach/.store in=\KAfach,         fach=,
  klasse/.store in=\KAklasse,     klasse=,
  thema/.store in=\KAthema,       thema=,
  datum/.store in=\KAdatum,       datum=,
  zeit/.store in=\KAzeit,         zeit=,
  hilfsmittel/.store in=\KAhilfsmittel, hilfsmittel=,
  gruppe/.store in=\KAgruppe,     gruppe=}

\newcommand\Kopf{%
  {\sffamily\bfseries\Large \KAtitel\ \KAfach\ \KAklasse
   \ifx\KAgruppe\empty\else\hfill Gruppe \KAgruppe\fi}\par\smallskip
  {\sffamily\large \KAthema}\par\medskip
  \begin{tabular}{@{}p{0.52\linewidth}p{0.42\linewidth}@{}}
    Name: \hrulefill & Datum: \KAdatum \\[4pt]
    Arbeitszeit: \KAzeit & Hilfsmittel: \KAhilfsmittel \\[4pt]
    Erreichbar: \KAzahl{\BEsumme}\,BE & Erreicht: \rule{1.6cm}{0.4pt}\,BE
      \quad Note: \rule{1.2cm}{0.4pt}
  \end{tabular}\par
  \smallskip\hrule\medskip}

% --- Punktetabelle ---------------------------------------------------
\newcommand\Punktetabelle{%
  \gdef\PTkopf{}\gdef\PTbe{}\gdef\PTleer{}%
  \ifnum\BEanzahl>0
  \foreach \i in {1,...,\BEanzahl}{%
    \xappto\PTkopf{& (\i)}%
    \xappto\PTbe{& \noexpand\KAzahl{\csuse{BEaufg\i}}}%
    \gappto\PTleer{&}}%
  \fi
  \begin{tabular}{@{}l*{\BEanzahl}{>{\centering\arraybackslash}p{1cm}}|>{\centering\arraybackslash}p{1.3cm}@{}}
    Aufgabe \PTkopf & $\Sigma$ \\ \hline
    BE \PTbe & \KAzahl{\BEsumme} \\[3pt] \hline
    erreicht \PTleer & \rule{0pt}{16pt}\\ \hline
  \end{tabular}}

% --- Notenschlüssel ----------------------------------------------------
% \Grenzen{95,80,60,40,20}: Mindestprozent für die Noten 1 bis 5.
% Die BE-Grenze wird auf halbe BE AUFgerundet.
\newcommand\Grenzen[1]{\def\KAgrenzen{#1}}
\Grenzen{95,80,60,40,20}
\newcommand\abBE[1]{\fpeval{ceil(#1/100*\BEsumme*2)/2}}
\newcommand\Notenschluessel{%
  \edef\@g{\KAgrenzen}%
  \expandafter\KA@noten\@g\relax}
\def\KA@noten#1,#2,#3,#4,#5\relax{%
  \begingroup\small
  \begin{tabular}{@{}l*{6}{>{\centering\arraybackslash}p{1.45cm}}@{}}
    Note & 1 & 2 & 3 & 4 & 5 & 6 \\ \hline
    ab \% & \KAzahl{#1} & \KAzahl{#2} & \KAzahl{#3} & \KAzahl{#4} & \KAzahl{#5} & 0 \\
    BE & \KA@bereich{\abBE{#1}}{\BEsumme}
       & \KA@bereich{\abBE{#2}}{\fpeval{\abBE{#1}-0.5}}
       & \KA@bereich{\abBE{#3}}{\fpeval{\abBE{#2}-0.5}}
       & \KA@bereich{\abBE{#4}}{\fpeval{\abBE{#3}-0.5}}
       & \KA@bereich{\abBE{#5}}{\fpeval{\abBE{#4}-0.5}}
       & \KA@bereich{0}{\fpeval{\abBE{#5}-0.5}} \\[2pt] \hline
    Anzahl & & & & & & \rule{0pt}{14pt}\\
  \end{tabular}\endgroup}
\def\KA@bereich#1#2{\KAzahl{#1}\,--\,\KAzahl{#2}}

% --- Lösungsteil -------------------------------------------------------
% Lösungen parallel zum Aufgabenteil, \scriptsize, Ergebnisse in \boxed{}
\newenvironment{loesungsteil}{%
  \clearpage\scriptsize\raggedright
  {\sffamily\bfseries\normalsize Lösungen \KAtitel\ \KAfach\ \KAklasse}\par\smallskip
  \hrule\medskip}{}
\makeatother
 nach \documentclass – oder nach
%  \input{header.tex} und \input{newcommands.tex} (universaler Stil).
%
%  \BE{2}            Bewertungseinheiten an eine (Teil-)Aufgabe hängen,
%                    halbe BE gehen auch: \BE{1.5}
%  \Kopf             Name/Klasse/Datum-Zeile + Angaben aus \Arbeit{...}
%  \Punktetabelle    Tabelle Aufgabe | BE | erreicht  (automatisch)
%  \Notenschluessel  Notentabelle nach \Grenzen (automatisch)
%  \loesungtrue      Lösungsteil mit ausgeben (\loesungfalse: ohne)
%
%  Alle Summen stehen erst nach dem ZWEITEN pdflatex-Lauf richtig da
%  (beim ersten Lauf erscheinen ??).
% ===================================================================
\makeatletter
% Pakete nur laden, wenn sie fehlen: so funktioniert die Datei allein
% (\documentclass + % ===================================================================
%  ka_makros.tex – Makros für Klassenarbeiten / Leistungskontrollen
%  Einbinden mit % ===================================================================
%  ka_makros.tex – Makros für Klassenarbeiten / Leistungskontrollen
%  Einbinden mit \input{ka_makros} nach \documentclass – oder nach
%  \input{header.tex} und \input{newcommands.tex} (universaler Stil).
%
%  \BE{2}            Bewertungseinheiten an eine (Teil-)Aufgabe hängen,
%                    halbe BE gehen auch: \BE{1.5}
%  \Kopf             Name/Klasse/Datum-Zeile + Angaben aus \Arbeit{...}
%  \Punktetabelle    Tabelle Aufgabe | BE | erreicht  (automatisch)
%  \Notenschluessel  Notentabelle nach \Grenzen (automatisch)
%  \loesungtrue      Lösungsteil mit ausgeben (\loesungfalse: ohne)
%
%  Alle Summen stehen erst nach dem ZWEITEN pdflatex-Lauf richtig da
%  (beim ersten Lauf erscheinen ??).
% ===================================================================
\makeatletter
% Pakete nur laden, wenn sie fehlen: so funktioniert die Datei allein
% (\documentclass + \input{ka_makros}) und nach der universalen header.tex
\@ifpackageloaded{fontenc}{}{\usepackage[T1]{fontenc}}
\@ifpackageloaded{lmodern}{}{\usepackage{lmodern}}
\@ifpackageloaded{babel}{}{\usepackage[ngerman]{babel}}
\@ifpackageloaded{geometry}{}{\usepackage[a4paper,margin=2cm]{geometry}}
\@ifpackageloaded{amsmath}{}{\usepackage{amsmath,amssymb}}
\@ifpackageloaded{xcolor}{}{\usepackage{xcolor}}
\@ifpackageloaded{enumitem}{}{\usepackage{enumitem}}
\@ifpackageloaded{multicol}{}{\usepackage{multicol}}
\@ifpackageloaded{booktabs}{}{\usepackage{array,booktabs}}
\@ifpackageloaded{etoolbox}{}{\usepackage{etoolbox}}
\@ifpackageloaded{tikz}{}{\usepackage{tikz}}
% Kein siunitx (kollidiert mit units aus der universalen header.tex):
% Zahlen mit Dezimalkomma über \KAzahl{1.5} -> 1,5
\newcommand\KAzahl[1]{\pgfmathprintnumber[use comma]{#1}}
\pagestyle{empty}                 % keine Kopf-/Fußzeilen
\setlength{\parindent}{0pt}


% --- Aufgabennummerierung (1), (2) … mit (a), (b) … ------------------
\setlist[enumerate,1]{label=(\arabic*), leftmargin=2.2em, itemsep=8pt}
\setlist[enumerate,2]{label=(\alph*), leftmargin=1.8em, itemsep=3pt}

% --- Schwierigkeitsstufen (falls nicht schon in newcommands.tex) ----
\providecommand\stufe[1]{%
  \ifcase#1\or
    \rule{1.3pt}{0.8ex}\,\textcolor{lightgray}{\rule{1.3pt}{1.3ex}\,\rule{1.3pt}{1.8ex}}\or
    \rule{1.3pt}{0.8ex}\,\rule{1.3pt}{1.3ex}\,\textcolor{lightgray}{\rule{1.3pt}{1.8ex}}\or
    \rule{1.3pt}{0.8ex}\,\rule{1.3pt}{1.3ex}\,\rule{1.3pt}{1.8ex}\fi\enspace}
\providecommand\leicht{\stufe{1}}
\providecommand\mittel{\stufe{2}}
\providecommand\schwer{\stufe{3}}

% --- Lösungsschalter --------------------------------------------------
\newif\ifloesung

% --- Bewertungseinheiten zählen ---------------------------------------
\gdef\BE@gesamt{0}
\gdef\BE@max{0}
\newcommand\BE[1]{%
  \edef\BE@nr{\arabic{enumi}}%
  \ifcsundef{BE@a\BE@nr}{\csgdef{BE@a\BE@nr}{0}}{}%
  \csxdef{BE@a\BE@nr}{\fpeval{\csuse{BE@a\BE@nr}+#1}}%
  \xdef\BE@gesamt{\fpeval{\BE@gesamt+#1}}%
  \ifnum\BE@nr>\BE@max\relax\xdef\BE@max{\BE@nr}\fi
  \unskip\nobreak\hfill\mbox{\enspace\textcolor{black!60}{\small[\KAzahl{#1}\,BE]}}}
% Am Ende in die .aux schreiben, damit der Kopf die Summe kennt
\AtEndDocument{%
  \immediate\write\@auxout{\string\gdef\string\BEsumme{\BE@gesamt}}%
  \immediate\write\@auxout{\string\gdef\string\BEanzahl{\BE@max}}%
  \@tempcnta=1\relax
  \loop\ifnum\@tempcnta>\BE@max\relax\else
    \immediate\write\@auxout{\string\csgdef{BEaufg\the\@tempcnta}{%
      \ifcsdef{BE@a\the\@tempcnta}{\csuse{BE@a\the\@tempcnta}}{0}}}%
    \advance\@tempcnta 1\relax\repeat}
\providecommand\BEsumme{0}  % erster Lauf: 0, ab dem zweiten Lauf korrekt
\providecommand\BEanzahl{0}

% --- Angaben zur Arbeit ------------------------------------------------
\newcommand\Arbeit[1]{\pgfkeys{/ka/.cd,#1}}
\pgfkeys{/ka/.cd,
  titel/.store in=\KAtitel,       titel=Klassenarbeit,
  fach/.store in=\KAfach,         fach=,
  klasse/.store in=\KAklasse,     klasse=,
  thema/.store in=\KAthema,       thema=,
  datum/.store in=\KAdatum,       datum=,
  zeit/.store in=\KAzeit,         zeit=,
  hilfsmittel/.store in=\KAhilfsmittel, hilfsmittel=,
  gruppe/.store in=\KAgruppe,     gruppe=}

\newcommand\Kopf{%
  {\sffamily\bfseries\Large \KAtitel\ \KAfach\ \KAklasse
   \ifx\KAgruppe\empty\else\hfill Gruppe \KAgruppe\fi}\par\smallskip
  {\sffamily\large \KAthema}\par\medskip
  \begin{tabular}{@{}p{0.52\linewidth}p{0.42\linewidth}@{}}
    Name: \hrulefill & Datum: \KAdatum \\[4pt]
    Arbeitszeit: \KAzeit & Hilfsmittel: \KAhilfsmittel \\[4pt]
    Erreichbar: \KAzahl{\BEsumme}\,BE & Erreicht: \rule{1.6cm}{0.4pt}\,BE
      \quad Note: \rule{1.2cm}{0.4pt}
  \end{tabular}\par
  \smallskip\hrule\medskip}

% --- Punktetabelle ---------------------------------------------------
\newcommand\Punktetabelle{%
  \gdef\PTkopf{}\gdef\PTbe{}\gdef\PTleer{}%
  \ifnum\BEanzahl>0
  \foreach \i in {1,...,\BEanzahl}{%
    \xappto\PTkopf{& (\i)}%
    \xappto\PTbe{& \noexpand\KAzahl{\csuse{BEaufg\i}}}%
    \gappto\PTleer{&}}%
  \fi
  \begin{tabular}{@{}l*{\BEanzahl}{>{\centering\arraybackslash}p{1cm}}|>{\centering\arraybackslash}p{1.3cm}@{}}
    Aufgabe \PTkopf & $\Sigma$ \\ \hline
    BE \PTbe & \KAzahl{\BEsumme} \\[3pt] \hline
    erreicht \PTleer & \rule{0pt}{16pt}\\ \hline
  \end{tabular}}

% --- Notenschlüssel ----------------------------------------------------
% \Grenzen{95,80,60,40,20}: Mindestprozent für die Noten 1 bis 5.
% Die BE-Grenze wird auf halbe BE AUFgerundet.
\newcommand\Grenzen[1]{\def\KAgrenzen{#1}}
\Grenzen{95,80,60,40,20}
\newcommand\abBE[1]{\fpeval{ceil(#1/100*\BEsumme*2)/2}}
\newcommand\Notenschluessel{%
  \edef\@g{\KAgrenzen}%
  \expandafter\KA@noten\@g\relax}
\def\KA@noten#1,#2,#3,#4,#5\relax{%
  \begingroup\small
  \begin{tabular}{@{}l*{6}{>{\centering\arraybackslash}p{1.45cm}}@{}}
    Note & 1 & 2 & 3 & 4 & 5 & 6 \\ \hline
    ab \% & \KAzahl{#1} & \KAzahl{#2} & \KAzahl{#3} & \KAzahl{#4} & \KAzahl{#5} & 0 \\
    BE & \KA@bereich{\abBE{#1}}{\BEsumme}
       & \KA@bereich{\abBE{#2}}{\fpeval{\abBE{#1}-0.5}}
       & \KA@bereich{\abBE{#3}}{\fpeval{\abBE{#2}-0.5}}
       & \KA@bereich{\abBE{#4}}{\fpeval{\abBE{#3}-0.5}}
       & \KA@bereich{\abBE{#5}}{\fpeval{\abBE{#4}-0.5}}
       & \KA@bereich{0}{\fpeval{\abBE{#5}-0.5}} \\[2pt] \hline
    Anzahl & & & & & & \rule{0pt}{14pt}\\
  \end{tabular}\endgroup}
\def\KA@bereich#1#2{\KAzahl{#1}\,--\,\KAzahl{#2}}

% --- Lösungsteil -------------------------------------------------------
% Lösungen parallel zum Aufgabenteil, \scriptsize, Ergebnisse in \boxed{}
\newenvironment{loesungsteil}{%
  \clearpage\scriptsize\raggedright
  {\sffamily\bfseries\normalsize Lösungen \KAtitel\ \KAfach\ \KAklasse}\par\smallskip
  \hrule\medskip}{}
\makeatother
 nach \documentclass – oder nach
%  \input{header.tex} und \input{newcommands.tex} (universaler Stil).
%
%  \BE{2}            Bewertungseinheiten an eine (Teil-)Aufgabe hängen,
%                    halbe BE gehen auch: \BE{1.5}
%  \Kopf             Name/Klasse/Datum-Zeile + Angaben aus \Arbeit{...}
%  \Punktetabelle    Tabelle Aufgabe | BE | erreicht  (automatisch)
%  \Notenschluessel  Notentabelle nach \Grenzen (automatisch)
%  \loesungtrue      Lösungsteil mit ausgeben (\loesungfalse: ohne)
%
%  Alle Summen stehen erst nach dem ZWEITEN pdflatex-Lauf richtig da
%  (beim ersten Lauf erscheinen ??).
% ===================================================================
\makeatletter
% Pakete nur laden, wenn sie fehlen: so funktioniert die Datei allein
% (\documentclass + % ===================================================================
%  ka_makros.tex – Makros für Klassenarbeiten / Leistungskontrollen
%  Einbinden mit \input{ka_makros} nach \documentclass – oder nach
%  \input{header.tex} und \input{newcommands.tex} (universaler Stil).
%
%  \BE{2}            Bewertungseinheiten an eine (Teil-)Aufgabe hängen,
%                    halbe BE gehen auch: \BE{1.5}
%  \Kopf             Name/Klasse/Datum-Zeile + Angaben aus \Arbeit{...}
%  \Punktetabelle    Tabelle Aufgabe | BE | erreicht  (automatisch)
%  \Notenschluessel  Notentabelle nach \Grenzen (automatisch)
%  \loesungtrue      Lösungsteil mit ausgeben (\loesungfalse: ohne)
%
%  Alle Summen stehen erst nach dem ZWEITEN pdflatex-Lauf richtig da
%  (beim ersten Lauf erscheinen ??).
% ===================================================================
\makeatletter
% Pakete nur laden, wenn sie fehlen: so funktioniert die Datei allein
% (\documentclass + \input{ka_makros}) und nach der universalen header.tex
\@ifpackageloaded{fontenc}{}{\usepackage[T1]{fontenc}}
\@ifpackageloaded{lmodern}{}{\usepackage{lmodern}}
\@ifpackageloaded{babel}{}{\usepackage[ngerman]{babel}}
\@ifpackageloaded{geometry}{}{\usepackage[a4paper,margin=2cm]{geometry}}
\@ifpackageloaded{amsmath}{}{\usepackage{amsmath,amssymb}}
\@ifpackageloaded{xcolor}{}{\usepackage{xcolor}}
\@ifpackageloaded{enumitem}{}{\usepackage{enumitem}}
\@ifpackageloaded{multicol}{}{\usepackage{multicol}}
\@ifpackageloaded{booktabs}{}{\usepackage{array,booktabs}}
\@ifpackageloaded{etoolbox}{}{\usepackage{etoolbox}}
\@ifpackageloaded{tikz}{}{\usepackage{tikz}}
% Kein siunitx (kollidiert mit units aus der universalen header.tex):
% Zahlen mit Dezimalkomma über \KAzahl{1.5} -> 1,5
\newcommand\KAzahl[1]{\pgfmathprintnumber[use comma]{#1}}
\pagestyle{empty}                 % keine Kopf-/Fußzeilen
\setlength{\parindent}{0pt}


% --- Aufgabennummerierung (1), (2) … mit (a), (b) … ------------------
\setlist[enumerate,1]{label=(\arabic*), leftmargin=2.2em, itemsep=8pt}
\setlist[enumerate,2]{label=(\alph*), leftmargin=1.8em, itemsep=3pt}

% --- Schwierigkeitsstufen (falls nicht schon in newcommands.tex) ----
\providecommand\stufe[1]{%
  \ifcase#1\or
    \rule{1.3pt}{0.8ex}\,\textcolor{lightgray}{\rule{1.3pt}{1.3ex}\,\rule{1.3pt}{1.8ex}}\or
    \rule{1.3pt}{0.8ex}\,\rule{1.3pt}{1.3ex}\,\textcolor{lightgray}{\rule{1.3pt}{1.8ex}}\or
    \rule{1.3pt}{0.8ex}\,\rule{1.3pt}{1.3ex}\,\rule{1.3pt}{1.8ex}\fi\enspace}
\providecommand\leicht{\stufe{1}}
\providecommand\mittel{\stufe{2}}
\providecommand\schwer{\stufe{3}}

% --- Lösungsschalter --------------------------------------------------
\newif\ifloesung

% --- Bewertungseinheiten zählen ---------------------------------------
\gdef\BE@gesamt{0}
\gdef\BE@max{0}
\newcommand\BE[1]{%
  \edef\BE@nr{\arabic{enumi}}%
  \ifcsundef{BE@a\BE@nr}{\csgdef{BE@a\BE@nr}{0}}{}%
  \csxdef{BE@a\BE@nr}{\fpeval{\csuse{BE@a\BE@nr}+#1}}%
  \xdef\BE@gesamt{\fpeval{\BE@gesamt+#1}}%
  \ifnum\BE@nr>\BE@max\relax\xdef\BE@max{\BE@nr}\fi
  \unskip\nobreak\hfill\mbox{\enspace\textcolor{black!60}{\small[\KAzahl{#1}\,BE]}}}
% Am Ende in die .aux schreiben, damit der Kopf die Summe kennt
\AtEndDocument{%
  \immediate\write\@auxout{\string\gdef\string\BEsumme{\BE@gesamt}}%
  \immediate\write\@auxout{\string\gdef\string\BEanzahl{\BE@max}}%
  \@tempcnta=1\relax
  \loop\ifnum\@tempcnta>\BE@max\relax\else
    \immediate\write\@auxout{\string\csgdef{BEaufg\the\@tempcnta}{%
      \ifcsdef{BE@a\the\@tempcnta}{\csuse{BE@a\the\@tempcnta}}{0}}}%
    \advance\@tempcnta 1\relax\repeat}
\providecommand\BEsumme{0}  % erster Lauf: 0, ab dem zweiten Lauf korrekt
\providecommand\BEanzahl{0}

% --- Angaben zur Arbeit ------------------------------------------------
\newcommand\Arbeit[1]{\pgfkeys{/ka/.cd,#1}}
\pgfkeys{/ka/.cd,
  titel/.store in=\KAtitel,       titel=Klassenarbeit,
  fach/.store in=\KAfach,         fach=,
  klasse/.store in=\KAklasse,     klasse=,
  thema/.store in=\KAthema,       thema=,
  datum/.store in=\KAdatum,       datum=,
  zeit/.store in=\KAzeit,         zeit=,
  hilfsmittel/.store in=\KAhilfsmittel, hilfsmittel=,
  gruppe/.store in=\KAgruppe,     gruppe=}

\newcommand\Kopf{%
  {\sffamily\bfseries\Large \KAtitel\ \KAfach\ \KAklasse
   \ifx\KAgruppe\empty\else\hfill Gruppe \KAgruppe\fi}\par\smallskip
  {\sffamily\large \KAthema}\par\medskip
  \begin{tabular}{@{}p{0.52\linewidth}p{0.42\linewidth}@{}}
    Name: \hrulefill & Datum: \KAdatum \\[4pt]
    Arbeitszeit: \KAzeit & Hilfsmittel: \KAhilfsmittel \\[4pt]
    Erreichbar: \KAzahl{\BEsumme}\,BE & Erreicht: \rule{1.6cm}{0.4pt}\,BE
      \quad Note: \rule{1.2cm}{0.4pt}
  \end{tabular}\par
  \smallskip\hrule\medskip}

% --- Punktetabelle ---------------------------------------------------
\newcommand\Punktetabelle{%
  \gdef\PTkopf{}\gdef\PTbe{}\gdef\PTleer{}%
  \ifnum\BEanzahl>0
  \foreach \i in {1,...,\BEanzahl}{%
    \xappto\PTkopf{& (\i)}%
    \xappto\PTbe{& \noexpand\KAzahl{\csuse{BEaufg\i}}}%
    \gappto\PTleer{&}}%
  \fi
  \begin{tabular}{@{}l*{\BEanzahl}{>{\centering\arraybackslash}p{1cm}}|>{\centering\arraybackslash}p{1.3cm}@{}}
    Aufgabe \PTkopf & $\Sigma$ \\ \hline
    BE \PTbe & \KAzahl{\BEsumme} \\[3pt] \hline
    erreicht \PTleer & \rule{0pt}{16pt}\\ \hline
  \end{tabular}}

% --- Notenschlüssel ----------------------------------------------------
% \Grenzen{95,80,60,40,20}: Mindestprozent für die Noten 1 bis 5.
% Die BE-Grenze wird auf halbe BE AUFgerundet.
\newcommand\Grenzen[1]{\def\KAgrenzen{#1}}
\Grenzen{95,80,60,40,20}
\newcommand\abBE[1]{\fpeval{ceil(#1/100*\BEsumme*2)/2}}
\newcommand\Notenschluessel{%
  \edef\@g{\KAgrenzen}%
  \expandafter\KA@noten\@g\relax}
\def\KA@noten#1,#2,#3,#4,#5\relax{%
  \begingroup\small
  \begin{tabular}{@{}l*{6}{>{\centering\arraybackslash}p{1.45cm}}@{}}
    Note & 1 & 2 & 3 & 4 & 5 & 6 \\ \hline
    ab \% & \KAzahl{#1} & \KAzahl{#2} & \KAzahl{#3} & \KAzahl{#4} & \KAzahl{#5} & 0 \\
    BE & \KA@bereich{\abBE{#1}}{\BEsumme}
       & \KA@bereich{\abBE{#2}}{\fpeval{\abBE{#1}-0.5}}
       & \KA@bereich{\abBE{#3}}{\fpeval{\abBE{#2}-0.5}}
       & \KA@bereich{\abBE{#4}}{\fpeval{\abBE{#3}-0.5}}
       & \KA@bereich{\abBE{#5}}{\fpeval{\abBE{#4}-0.5}}
       & \KA@bereich{0}{\fpeval{\abBE{#5}-0.5}} \\[2pt] \hline
    Anzahl & & & & & & \rule{0pt}{14pt}\\
  \end{tabular}\endgroup}
\def\KA@bereich#1#2{\KAzahl{#1}\,--\,\KAzahl{#2}}

% --- Lösungsteil -------------------------------------------------------
% Lösungen parallel zum Aufgabenteil, \scriptsize, Ergebnisse in \boxed{}
\newenvironment{loesungsteil}{%
  \clearpage\scriptsize\raggedright
  {\sffamily\bfseries\normalsize Lösungen \KAtitel\ \KAfach\ \KAklasse}\par\smallskip
  \hrule\medskip}{}
\makeatother
) und nach der universalen header.tex
\@ifpackageloaded{fontenc}{}{\usepackage[T1]{fontenc}}
\@ifpackageloaded{lmodern}{}{\usepackage{lmodern}}
\@ifpackageloaded{babel}{}{\usepackage[ngerman]{babel}}
\@ifpackageloaded{geometry}{}{\usepackage[a4paper,margin=2cm]{geometry}}
\@ifpackageloaded{amsmath}{}{\usepackage{amsmath,amssymb}}
\@ifpackageloaded{xcolor}{}{\usepackage{xcolor}}
\@ifpackageloaded{enumitem}{}{\usepackage{enumitem}}
\@ifpackageloaded{multicol}{}{\usepackage{multicol}}
\@ifpackageloaded{booktabs}{}{\usepackage{array,booktabs}}
\@ifpackageloaded{etoolbox}{}{\usepackage{etoolbox}}
\@ifpackageloaded{tikz}{}{\usepackage{tikz}}
% Kein siunitx (kollidiert mit units aus der universalen header.tex):
% Zahlen mit Dezimalkomma über \KAzahl{1.5} -> 1,5
\newcommand\KAzahl[1]{\pgfmathprintnumber[use comma]{#1}}
\pagestyle{empty}                 % keine Kopf-/Fußzeilen
\setlength{\parindent}{0pt}


% --- Aufgabennummerierung (1), (2) … mit (a), (b) … ------------------
\setlist[enumerate,1]{label=(\arabic*), leftmargin=2.2em, itemsep=8pt}
\setlist[enumerate,2]{label=(\alph*), leftmargin=1.8em, itemsep=3pt}

% --- Schwierigkeitsstufen (falls nicht schon in newcommands.tex) ----
\providecommand\stufe[1]{%
  \ifcase#1\or
    \rule{1.3pt}{0.8ex}\,\textcolor{lightgray}{\rule{1.3pt}{1.3ex}\,\rule{1.3pt}{1.8ex}}\or
    \rule{1.3pt}{0.8ex}\,\rule{1.3pt}{1.3ex}\,\textcolor{lightgray}{\rule{1.3pt}{1.8ex}}\or
    \rule{1.3pt}{0.8ex}\,\rule{1.3pt}{1.3ex}\,\rule{1.3pt}{1.8ex}\fi\enspace}
\providecommand\leicht{\stufe{1}}
\providecommand\mittel{\stufe{2}}
\providecommand\schwer{\stufe{3}}

% --- Lösungsschalter --------------------------------------------------
\newif\ifloesung

% --- Bewertungseinheiten zählen ---------------------------------------
\gdef\BE@gesamt{0}
\gdef\BE@max{0}
\newcommand\BE[1]{%
  \edef\BE@nr{\arabic{enumi}}%
  \ifcsundef{BE@a\BE@nr}{\csgdef{BE@a\BE@nr}{0}}{}%
  \csxdef{BE@a\BE@nr}{\fpeval{\csuse{BE@a\BE@nr}+#1}}%
  \xdef\BE@gesamt{\fpeval{\BE@gesamt+#1}}%
  \ifnum\BE@nr>\BE@max\relax\xdef\BE@max{\BE@nr}\fi
  \unskip\nobreak\hfill\mbox{\enspace\textcolor{black!60}{\small[\KAzahl{#1}\,BE]}}}
% Am Ende in die .aux schreiben, damit der Kopf die Summe kennt
\AtEndDocument{%
  \immediate\write\@auxout{\string\gdef\string\BEsumme{\BE@gesamt}}%
  \immediate\write\@auxout{\string\gdef\string\BEanzahl{\BE@max}}%
  \@tempcnta=1\relax
  \loop\ifnum\@tempcnta>\BE@max\relax\else
    \immediate\write\@auxout{\string\csgdef{BEaufg\the\@tempcnta}{%
      \ifcsdef{BE@a\the\@tempcnta}{\csuse{BE@a\the\@tempcnta}}{0}}}%
    \advance\@tempcnta 1\relax\repeat}
\providecommand\BEsumme{0}  % erster Lauf: 0, ab dem zweiten Lauf korrekt
\providecommand\BEanzahl{0}

% --- Angaben zur Arbeit ------------------------------------------------
\newcommand\Arbeit[1]{\pgfkeys{/ka/.cd,#1}}
\pgfkeys{/ka/.cd,
  titel/.store in=\KAtitel,       titel=Klassenarbeit,
  fach/.store in=\KAfach,         fach=,
  klasse/.store in=\KAklasse,     klasse=,
  thema/.store in=\KAthema,       thema=,
  datum/.store in=\KAdatum,       datum=,
  zeit/.store in=\KAzeit,         zeit=,
  hilfsmittel/.store in=\KAhilfsmittel, hilfsmittel=,
  gruppe/.store in=\KAgruppe,     gruppe=}

\newcommand\Kopf{%
  {\sffamily\bfseries\Large \KAtitel\ \KAfach\ \KAklasse
   \ifx\KAgruppe\empty\else\hfill Gruppe \KAgruppe\fi}\par\smallskip
  {\sffamily\large \KAthema}\par\medskip
  \begin{tabular}{@{}p{0.52\linewidth}p{0.42\linewidth}@{}}
    Name: \hrulefill & Datum: \KAdatum \\[4pt]
    Arbeitszeit: \KAzeit & Hilfsmittel: \KAhilfsmittel \\[4pt]
    Erreichbar: \KAzahl{\BEsumme}\,BE & Erreicht: \rule{1.6cm}{0.4pt}\,BE
      \quad Note: \rule{1.2cm}{0.4pt}
  \end{tabular}\par
  \smallskip\hrule\medskip}

% --- Punktetabelle ---------------------------------------------------
\newcommand\Punktetabelle{%
  \gdef\PTkopf{}\gdef\PTbe{}\gdef\PTleer{}%
  \ifnum\BEanzahl>0
  \foreach \i in {1,...,\BEanzahl}{%
    \xappto\PTkopf{& (\i)}%
    \xappto\PTbe{& \noexpand\KAzahl{\csuse{BEaufg\i}}}%
    \gappto\PTleer{&}}%
  \fi
  \begin{tabular}{@{}l*{\BEanzahl}{>{\centering\arraybackslash}p{1cm}}|>{\centering\arraybackslash}p{1.3cm}@{}}
    Aufgabe \PTkopf & $\Sigma$ \\ \hline
    BE \PTbe & \KAzahl{\BEsumme} \\[3pt] \hline
    erreicht \PTleer & \rule{0pt}{16pt}\\ \hline
  \end{tabular}}

% --- Notenschlüssel ----------------------------------------------------
% \Grenzen{95,80,60,40,20}: Mindestprozent für die Noten 1 bis 5.
% Die BE-Grenze wird auf halbe BE AUFgerundet.
\newcommand\Grenzen[1]{\def\KAgrenzen{#1}}
\Grenzen{95,80,60,40,20}
\newcommand\abBE[1]{\fpeval{ceil(#1/100*\BEsumme*2)/2}}
\newcommand\Notenschluessel{%
  \edef\@g{\KAgrenzen}%
  \expandafter\KA@noten\@g\relax}
\def\KA@noten#1,#2,#3,#4,#5\relax{%
  \begingroup\small
  \begin{tabular}{@{}l*{6}{>{\centering\arraybackslash}p{1.45cm}}@{}}
    Note & 1 & 2 & 3 & 4 & 5 & 6 \\ \hline
    ab \% & \KAzahl{#1} & \KAzahl{#2} & \KAzahl{#3} & \KAzahl{#4} & \KAzahl{#5} & 0 \\
    BE & \KA@bereich{\abBE{#1}}{\BEsumme}
       & \KA@bereich{\abBE{#2}}{\fpeval{\abBE{#1}-0.5}}
       & \KA@bereich{\abBE{#3}}{\fpeval{\abBE{#2}-0.5}}
       & \KA@bereich{\abBE{#4}}{\fpeval{\abBE{#3}-0.5}}
       & \KA@bereich{\abBE{#5}}{\fpeval{\abBE{#4}-0.5}}
       & \KA@bereich{0}{\fpeval{\abBE{#5}-0.5}} \\[2pt] \hline
    Anzahl & & & & & & \rule{0pt}{14pt}\\
  \end{tabular}\endgroup}
\def\KA@bereich#1#2{\KAzahl{#1}\,--\,\KAzahl{#2}}

% --- Lösungsteil -------------------------------------------------------
% Lösungen parallel zum Aufgabenteil, \scriptsize, Ergebnisse in \boxed{}
\newenvironment{loesungsteil}{%
  \clearpage\scriptsize\raggedright
  {\sffamily\bfseries\normalsize Lösungen \KAtitel\ \KAfach\ \KAklasse}\par\smallskip
  \hrule\medskip}{}
\makeatother
) und nach der universalen header.tex
\@ifpackageloaded{fontenc}{}{\usepackage[T1]{fontenc}}
\@ifpackageloaded{lmodern}{}{\usepackage{lmodern}}
\@ifpackageloaded{babel}{}{\usepackage[ngerman]{babel}}
\@ifpackageloaded{geometry}{}{\usepackage[a4paper,margin=2cm]{geometry}}
\@ifpackageloaded{amsmath}{}{\usepackage{amsmath,amssymb}}
\@ifpackageloaded{xcolor}{}{\usepackage{xcolor}}
\@ifpackageloaded{enumitem}{}{\usepackage{enumitem}}
\@ifpackageloaded{multicol}{}{\usepackage{multicol}}
\@ifpackageloaded{booktabs}{}{\usepackage{array,booktabs}}
\@ifpackageloaded{etoolbox}{}{\usepackage{etoolbox}}
\@ifpackageloaded{tikz}{}{\usepackage{tikz}}
% Kein siunitx (kollidiert mit units aus der universalen header.tex):
% Zahlen mit Dezimalkomma über \KAzahl{1.5} -> 1,5
\newcommand\KAzahl[1]{\pgfmathprintnumber[use comma]{#1}}
\pagestyle{empty}                 % keine Kopf-/Fußzeilen
\setlength{\parindent}{0pt}


% --- Aufgabennummerierung (1), (2) … mit (a), (b) … ------------------
\setlist[enumerate,1]{label=(\arabic*), leftmargin=2.2em, itemsep=8pt}
\setlist[enumerate,2]{label=(\alph*), leftmargin=1.8em, itemsep=3pt}

% --- Schwierigkeitsstufen (falls nicht schon in newcommands.tex) ----
\providecommand\stufe[1]{%
  \ifcase#1\or
    \rule{1.3pt}{0.8ex}\,\textcolor{lightgray}{\rule{1.3pt}{1.3ex}\,\rule{1.3pt}{1.8ex}}\or
    \rule{1.3pt}{0.8ex}\,\rule{1.3pt}{1.3ex}\,\textcolor{lightgray}{\rule{1.3pt}{1.8ex}}\or
    \rule{1.3pt}{0.8ex}\,\rule{1.3pt}{1.3ex}\,\rule{1.3pt}{1.8ex}\fi\enspace}
\providecommand\leicht{\stufe{1}}
\providecommand\mittel{\stufe{2}}
\providecommand\schwer{\stufe{3}}

% --- Lösungsschalter --------------------------------------------------
\newif\ifloesung

% --- Bewertungseinheiten zählen ---------------------------------------
\gdef\BE@gesamt{0}
\gdef\BE@max{0}
\newcommand\BE[1]{%
  \edef\BE@nr{\arabic{enumi}}%
  \ifcsundef{BE@a\BE@nr}{\csgdef{BE@a\BE@nr}{0}}{}%
  \csxdef{BE@a\BE@nr}{\fpeval{\csuse{BE@a\BE@nr}+#1}}%
  \xdef\BE@gesamt{\fpeval{\BE@gesamt+#1}}%
  \ifnum\BE@nr>\BE@max\relax\xdef\BE@max{\BE@nr}\fi
  \unskip\nobreak\hfill\mbox{\enspace\textcolor{black!60}{\small[\KAzahl{#1}\,BE]}}}
% Am Ende in die .aux schreiben, damit der Kopf die Summe kennt
\AtEndDocument{%
  \immediate\write\@auxout{\string\gdef\string\BEsumme{\BE@gesamt}}%
  \immediate\write\@auxout{\string\gdef\string\BEanzahl{\BE@max}}%
  \@tempcnta=1\relax
  \loop\ifnum\@tempcnta>\BE@max\relax\else
    \immediate\write\@auxout{\string\csgdef{BEaufg\the\@tempcnta}{%
      \ifcsdef{BE@a\the\@tempcnta}{\csuse{BE@a\the\@tempcnta}}{0}}}%
    \advance\@tempcnta 1\relax\repeat}
\providecommand\BEsumme{0}  % erster Lauf: 0, ab dem zweiten Lauf korrekt
\providecommand\BEanzahl{0}

% --- Angaben zur Arbeit ------------------------------------------------
\newcommand\Arbeit[1]{\pgfkeys{/ka/.cd,#1}}
\pgfkeys{/ka/.cd,
  titel/.store in=\KAtitel,       titel=Klassenarbeit,
  fach/.store in=\KAfach,         fach=,
  klasse/.store in=\KAklasse,     klasse=,
  thema/.store in=\KAthema,       thema=,
  datum/.store in=\KAdatum,       datum=,
  zeit/.store in=\KAzeit,         zeit=,
  hilfsmittel/.store in=\KAhilfsmittel, hilfsmittel=,
  gruppe/.store in=\KAgruppe,     gruppe=}

\newcommand\Kopf{%
  {\sffamily\bfseries\Large \KAtitel\ \KAfach\ \KAklasse
   \ifx\KAgruppe\empty\else\hfill Gruppe \KAgruppe\fi}\par\smallskip
  {\sffamily\large \KAthema}\par\medskip
  \begin{tabular}{@{}p{0.52\linewidth}p{0.42\linewidth}@{}}
    Name: \hrulefill & Datum: \KAdatum \\[4pt]
    Arbeitszeit: \KAzeit & Hilfsmittel: \KAhilfsmittel \\[4pt]
    Erreichbar: \KAzahl{\BEsumme}\,BE & Erreicht: \rule{1.6cm}{0.4pt}\,BE
      \quad Note: \rule{1.2cm}{0.4pt}
  \end{tabular}\par
  \smallskip\hrule\medskip}

% --- Punktetabelle ---------------------------------------------------
\newcommand\Punktetabelle{%
  \gdef\PTkopf{}\gdef\PTbe{}\gdef\PTleer{}%
  \ifnum\BEanzahl>0
  \foreach \i in {1,...,\BEanzahl}{%
    \xappto\PTkopf{& (\i)}%
    \xappto\PTbe{& \noexpand\KAzahl{\csuse{BEaufg\i}}}%
    \gappto\PTleer{&}}%
  \fi
  \begin{tabular}{@{}l*{\BEanzahl}{>{\centering\arraybackslash}p{1cm}}|>{\centering\arraybackslash}p{1.3cm}@{}}
    Aufgabe \PTkopf & $\Sigma$ \\ \hline
    BE \PTbe & \KAzahl{\BEsumme} \\[3pt] \hline
    erreicht \PTleer & \rule{0pt}{16pt}\\ \hline
  \end{tabular}}

% --- Notenschlüssel ----------------------------------------------------
% \Grenzen{95,80,60,40,20}: Mindestprozent für die Noten 1 bis 5.
% Die BE-Grenze wird auf halbe BE AUFgerundet.
\newcommand\Grenzen[1]{\def\KAgrenzen{#1}}
\Grenzen{95,80,60,40,20}
\newcommand\abBE[1]{\fpeval{ceil(#1/100*\BEsumme*2)/2}}
\newcommand\Notenschluessel{%
  \edef\@g{\KAgrenzen}%
  \expandafter\KA@noten\@g\relax}
\def\KA@noten#1,#2,#3,#4,#5\relax{%
  \begingroup\small
  \begin{tabular}{@{}l*{6}{>{\centering\arraybackslash}p{1.45cm}}@{}}
    Note & 1 & 2 & 3 & 4 & 5 & 6 \\ \hline
    ab \% & \KAzahl{#1} & \KAzahl{#2} & \KAzahl{#3} & \KAzahl{#4} & \KAzahl{#5} & 0 \\
    BE & \KA@bereich{\abBE{#1}}{\BEsumme}
       & \KA@bereich{\abBE{#2}}{\fpeval{\abBE{#1}-0.5}}
       & \KA@bereich{\abBE{#3}}{\fpeval{\abBE{#2}-0.5}}
       & \KA@bereich{\abBE{#4}}{\fpeval{\abBE{#3}-0.5}}
       & \KA@bereich{\abBE{#5}}{\fpeval{\abBE{#4}-0.5}}
       & \KA@bereich{0}{\fpeval{\abBE{#5}-0.5}} \\[2pt] \hline
    Anzahl & & & & & & \rule{0pt}{14pt}\\
  \end{tabular}\endgroup}
\def\KA@bereich#1#2{\KAzahl{#1}\,--\,\KAzahl{#2}}

% --- Lösungsteil -------------------------------------------------------
% Lösungen parallel zum Aufgabenteil, \scriptsize, Ergebnisse in \boxed{}
\newenvironment{loesungsteil}{%
  \clearpage\scriptsize\raggedright
  {\sffamily\bfseries\normalsize Lösungen \KAtitel\ \KAfach\ \KAklasse}\par\smallskip
  \hrule\medskip}{}
\makeatother
) und nach der universalen header.tex
\@ifpackageloaded{fontenc}{}{\usepackage[T1]{fontenc}}
\@ifpackageloaded{lmodern}{}{\usepackage{lmodern}}
\@ifpackageloaded{babel}{}{\usepackage[ngerman]{babel}}
\@ifpackageloaded{geometry}{}{\usepackage[a4paper,margin=2cm]{geometry}}
\@ifpackageloaded{amsmath}{}{\usepackage{amsmath,amssymb}}
\@ifpackageloaded{xcolor}{}{\usepackage{xcolor}}
\@ifpackageloaded{enumitem}{}{\usepackage{enumitem}}
\@ifpackageloaded{multicol}{}{\usepackage{multicol}}
\@ifpackageloaded{booktabs}{}{\usepackage{array,booktabs}}
\@ifpackageloaded{etoolbox}{}{\usepackage{etoolbox}}
\@ifpackageloaded{tikz}{}{\usepackage{tikz}}
% Kein siunitx (kollidiert mit units aus der universalen header.tex):
% Zahlen mit Dezimalkomma über \KAzahl{1.5} -> 1,5
\newcommand\KAzahl[1]{\pgfmathprintnumber[use comma]{#1}}
\pagestyle{empty}                 % keine Kopf-/Fußzeilen
\setlength{\parindent}{0pt}


% --- Aufgabennummerierung (1), (2) … mit (a), (b) … ------------------
\setlist[enumerate,1]{label=(\arabic*), leftmargin=2.2em, itemsep=8pt}
\setlist[enumerate,2]{label=(\alph*), leftmargin=1.8em, itemsep=3pt}

% --- Schwierigkeitsstufen (falls nicht schon in newcommands.tex) ----
\providecommand\stufe[1]{%
  \ifcase#1\or
    \rule{1.3pt}{0.8ex}\,\textcolor{lightgray}{\rule{1.3pt}{1.3ex}\,\rule{1.3pt}{1.8ex}}\or
    \rule{1.3pt}{0.8ex}\,\rule{1.3pt}{1.3ex}\,\textcolor{lightgray}{\rule{1.3pt}{1.8ex}}\or
    \rule{1.3pt}{0.8ex}\,\rule{1.3pt}{1.3ex}\,\rule{1.3pt}{1.8ex}\fi\enspace}
\providecommand\leicht{\stufe{1}}
\providecommand\mittel{\stufe{2}}
\providecommand\schwer{\stufe{3}}

% --- Lösungsschalter --------------------------------------------------
\newif\ifloesung

% --- Bewertungseinheiten zählen ---------------------------------------
\gdef\BE@gesamt{0}
\gdef\BE@max{0}
\newcommand\BE[1]{%
  \edef\BE@nr{\arabic{enumi}}%
  \ifcsundef{BE@a\BE@nr}{\csgdef{BE@a\BE@nr}{0}}{}%
  \csxdef{BE@a\BE@nr}{\fpeval{\csuse{BE@a\BE@nr}+#1}}%
  \xdef\BE@gesamt{\fpeval{\BE@gesamt+#1}}%
  \ifnum\BE@nr>\BE@max\relax\xdef\BE@max{\BE@nr}\fi
  \unskip\nobreak\hfill\mbox{\enspace\textcolor{black!60}{\small[\KAzahl{#1}\,BE]}}}
% Am Ende in die .aux schreiben, damit der Kopf die Summe kennt
\AtEndDocument{%
  \immediate\write\@auxout{\string\gdef\string\BEsumme{\BE@gesamt}}%
  \immediate\write\@auxout{\string\gdef\string\BEanzahl{\BE@max}}%
  \@tempcnta=1\relax
  \loop\ifnum\@tempcnta>\BE@max\relax\else
    \immediate\write\@auxout{\string\csgdef{BEaufg\the\@tempcnta}{%
      \ifcsdef{BE@a\the\@tempcnta}{\csuse{BE@a\the\@tempcnta}}{0}}}%
    \advance\@tempcnta 1\relax\repeat}
\providecommand\BEsumme{0}  % erster Lauf: 0, ab dem zweiten Lauf korrekt
\providecommand\BEanzahl{0}

% --- Angaben zur Arbeit ------------------------------------------------
\newcommand\Arbeit[1]{\pgfkeys{/ka/.cd,#1}}
\pgfkeys{/ka/.cd,
  titel/.store in=\KAtitel,       titel=Klassenarbeit,
  fach/.store in=\KAfach,         fach=,
  klasse/.store in=\KAklasse,     klasse=,
  thema/.store in=\KAthema,       thema=,
  datum/.store in=\KAdatum,       datum=,
  zeit/.store in=\KAzeit,         zeit=,
  hilfsmittel/.store in=\KAhilfsmittel, hilfsmittel=,
  gruppe/.store in=\KAgruppe,     gruppe=}

\newcommand\Kopf{%
  {\sffamily\bfseries\Large \KAtitel\ \KAfach\ \KAklasse
   \ifx\KAgruppe\empty\else\hfill Gruppe \KAgruppe\fi}\par\smallskip
  {\sffamily\large \KAthema}\par\medskip
  \begin{tabular}{@{}p{0.52\linewidth}p{0.42\linewidth}@{}}
    Name: \hrulefill & Datum: \KAdatum \\[4pt]
    Arbeitszeit: \KAzeit & Hilfsmittel: \KAhilfsmittel \\[4pt]
    Erreichbar: \KAzahl{\BEsumme}\,BE & Erreicht: \rule{1.6cm}{0.4pt}\,BE
      \quad Note: \rule{1.2cm}{0.4pt}
  \end{tabular}\par
  \smallskip\hrule\medskip}

% --- Punktetabelle ---------------------------------------------------
\newcommand\Punktetabelle{%
  \gdef\PTkopf{}\gdef\PTbe{}\gdef\PTleer{}%
  \ifnum\BEanzahl>0
  \foreach \i in {1,...,\BEanzahl}{%
    \xappto\PTkopf{& (\i)}%
    \xappto\PTbe{& \noexpand\KAzahl{\csuse{BEaufg\i}}}%
    \gappto\PTleer{&}}%
  \fi
  \begin{tabular}{@{}l*{\BEanzahl}{>{\centering\arraybackslash}p{1cm}}|>{\centering\arraybackslash}p{1.3cm}@{}}
    Aufgabe \PTkopf & $\Sigma$ \\ \hline
    BE \PTbe & \KAzahl{\BEsumme} \\[3pt] \hline
    erreicht \PTleer & \rule{0pt}{16pt}\\ \hline
  \end{tabular}}

% --- Notenschlüssel ----------------------------------------------------
% \Grenzen{95,80,60,40,20}: Mindestprozent für die Noten 1 bis 5.
% Die BE-Grenze wird auf halbe BE AUFgerundet.
\newcommand\Grenzen[1]{\def\KAgrenzen{#1}}
\Grenzen{95,80,60,40,20}
\newcommand\abBE[1]{\fpeval{ceil(#1/100*\BEsumme*2)/2}}
\newcommand\Notenschluessel{%
  \edef\@g{\KAgrenzen}%
  \expandafter\KA@noten\@g\relax}
\def\KA@noten#1,#2,#3,#4,#5\relax{%
  \begingroup\small
  \begin{tabular}{@{}l*{6}{>{\centering\arraybackslash}p{1.45cm}}@{}}
    Note & 1 & 2 & 3 & 4 & 5 & 6 \\ \hline
    ab \% & \KAzahl{#1} & \KAzahl{#2} & \KAzahl{#3} & \KAzahl{#4} & \KAzahl{#5} & 0 \\
    BE & \KA@bereich{\abBE{#1}}{\BEsumme}
       & \KA@bereich{\abBE{#2}}{\fpeval{\abBE{#1}-0.5}}
       & \KA@bereich{\abBE{#3}}{\fpeval{\abBE{#2}-0.5}}
       & \KA@bereich{\abBE{#4}}{\fpeval{\abBE{#3}-0.5}}
       & \KA@bereich{\abBE{#5}}{\fpeval{\abBE{#4}-0.5}}
       & \KA@bereich{0}{\fpeval{\abBE{#5}-0.5}} \\[2pt] \hline
    Anzahl & & & & & & \rule{0pt}{14pt}\\
  \end{tabular}\endgroup}
\def\KA@bereich#1#2{\KAzahl{#1}\,--\,\KAzahl{#2}}

% --- Lösungsteil -------------------------------------------------------
% Lösungen parallel zum Aufgabenteil, \scriptsize, Ergebnisse in \boxed{}
\newenvironment{loesungsteil}{%
  \clearpage\scriptsize\raggedright
  {\sffamily\bfseries\normalsize Lösungen \KAtitel\ \KAfach\ \KAklasse}\par\smallskip
  \hrule\medskip}{}
\makeatother
